\chapter{Homological Algebra}\label{homological-algebra}
\chapterstatus{complete}

\section{Introduction}\label{homological-algebra:section-introduction}

Homological algebra measures the failure of exactness. Its basic objects are complexes and their
cohomology; its basic tools are the long exact sequence and derived functors. Everything in this project
that is called a cohomology theory (singular cohomology in Chapter~\ref{singular-homology}, sheaf
cohomology in Chapter~\ref{sheaves}, de Rham and Dolbeault cohomology in
Chapters~\ref{de-rham} and~\ref{complex-manifolds}) is built with the machinery of this chapter.

We work in abelian categories, but prove the diagram lemmas for modules; by the Freyd--Mitchell embedding
theorem they then hold in every abelian category (Remark~\ref{homological-algebra:remark-freyd-mitchell}).
References are \cite{weibel-homological}, \cite{cartan-eilenberg} and \cite{gelfand-manin}.

\noindent\textbf{Prerequisites.} Chapter~\ref{categories} (Categories) and Chapter~\ref{rings} (Rings and
Modules).

\section{Abelian categories}\label{homological-algebra:section-abelian}

\begin{definition}\label{homological-algebra:definition-abelian}
An additive category $\mathcal{A}$ (Definition~\ref{categories:definition-additive}) is \emph{abelian} if every
morphism has a kernel and a cokernel (Definition~\ref{categories:definition-biproduct}), every monomorphism is a
kernel of its cokernel, and every epimorphism is a cokernel of its kernel. For $f \colon X \to Y$ the \emph{image} is
$\im f = \ker(\operatorname{coker} f)$. A sequence $X \xrightarrow{f} Y \xrightarrow{g} Z$ with $gf = 0$ is \emph{exact} at $Y$ if the
induced morphism $\im f \to \ker g$ is an isomorphism. An additive functor between abelian categories is
\emph{left exact} if it sends exact sequences $0 \to X \to Y \to Z$ to exact sequences, \emph{right exact} if it sends
exact $X \to Y \to Z \to 0$ to exact sequences, and \emph{exact} if both.
\end{definition}

\begin{example}\label{homological-algebra:example-abelian}
For a ring $R$, the category $R\text{-}\cat{Mod}$ is abelian (Proposition~\ref{rings:proposition-module-basics}):
kernels and cokernels are the usual ones, an injective map is the kernel of the projection onto its cokernel, and a
surjective map is the cokernel of the inclusion of its kernel. In particular $\cat{Ab}$ is abelian. The category of
sheaves of abelian groups on a topological space is abelian (Chapter~\ref{sheaves}). By
Theorem~\ref{rings:theorem-tensor-hom}, $- \otimes_R N$ is right exact, and $\Hom_R(N, -)$ is left exact.
\end{example}

\begin{lemma}\label{homological-algebra:lemma-mono-epi-iso}
In an abelian category a morphism that is both a monomorphism and an epimorphism is an isomorphism.
\end{lemma}

\begin{proof}
Let $f \colon X \to Y$ be both. Its cokernel is $Y \to 0$, since $f$ is an epimorphism; as $f$ is a monomorphism, it is a
kernel of $Y \to 0$. But $\id_Y$ is also a kernel of $Y \to 0$, so $f$ is an isomorphism by uniqueness of kernels
(Lemma~\ref{categories:lemma-limit-unique}).
\end{proof}

\begin{counterexample}\label{homological-algebra:counterexample-not-abelian}
The category $\mathcal{F}$ of finitely generated free abelian groups is additive and has kernels
(Proposition~\ref{groups:proposition-subgroups-free}), but it is not abelian. The morphism $2 \colon \ZZ \to \ZZ$ is a
monomorphism, and it is an epimorphism in $\mathcal{F}$: if $g \circ 2 = 0$ for $g \colon \ZZ \to \ZZ^n$, then $2g(1) = 0$, so
$g = 0$. It is not an isomorphism, contradicting Lemma~\ref{homological-algebra:lemma-mono-epi-iso}.
\end{counterexample}

\begin{remark}\label{homological-algebra:remark-freyd-mitchell}
By the Freyd--Mitchell embedding theorem, every small abelian category admits a fully faithful exact functor into
$R\text{-}\cat{Mod}$ for some ring $R$ \cite[Theorem 1.6.1]{weibel-homological}. Every statement about a finite diagram
in an abelian category $\mathcal{A}$ that asserts exactness, or the existence of morphisms constructed from kernels,
cokernels and compositions, involves only a small abelian subcategory of $\mathcal{A}$; hence it suffices to prove it for
modules, where one may argue with elements (``diagram chasing''). We use this principle for
Lemmas~\ref{homological-algebra:lemma-snake} and~\ref{homological-algebra:lemma-five}, and wherever we prove a
statement about abelian categories by a chase.
\end{remark}

\section{Complexes}\label{homological-algebra:section-complexes}

\begin{definition}\label{homological-algebra:definition-complex}
A \emph{complex} (or cochain complex) $A^\bullet$ in an abelian category $\mathcal{A}$ is a family of objects $A^n$,
$n \in \ZZ$, with morphisms $d^n \colon A^n \to A^{n+1}$ such that $d^{n+1}d^n = 0$. Its \emph{cocycles}, \emph{coboundaries} and
\emph{cohomology} are $Z^n = \ker d^n$, $B^n = \im d^{n-1}$ and $H^n(A^\bullet) = Z^n/B^n$. A \emph{morphism of complexes}
$f \colon A^\bullet \to B^\bullet$ is a family $f^n \colon A^n \to B^n$ with $df = fd$; it induces $H^n(f)$. It is a
\emph{quasi-isomorphism} if all $H^n(f)$ are isomorphisms. A complex is \emph{acyclic} (or exact) if all $H^n$ vanish, and
\emph{bounded below} if $A^n = 0$ for $n \ll 0$. The complexes form an abelian category $\mathrm{Ch}(\mathcal{A})$, with
kernels and cokernels computed degreewise.
\end{definition}

\begin{notation}\label{homological-algebra:notation-chain-complexes}
A \emph{chain complex} $C_\bullet$ has differentials $d_n \colon C_n \to C_{n-1}$ and \emph{homology}
$H_n = \ker d_n/\im d_{n+1}$. Setting $C^n = C_{-n}$ turns it into a complex, with $H^n = H_{-n}$; all statements about
complexes apply to chain complexes in this way. For a complex $A^\bullet$ and $k \in \ZZ$, the \emph{shift} $A[k]$ has
$A[k]^n = A^{n+k}$ and differential $(-1)^kd$.
\end{notation}

\begin{definition}\label{homological-algebra:definition-homotopy}
Two morphisms $f, g \colon A^\bullet \to B^\bullet$ are \emph{homotopic}, $f \simeq g$, if there are morphisms
$h^n \colon A^n \to B^{n-1}$ with $f^n - g^n = d^{n-1}h^n + h^{n+1}d^n$. A morphism $f$ is a \emph{homotopy equivalence} if there
is $g$ with $gf \simeq \id$ and $fg \simeq \id$.
\end{definition}

\begin{lemma}\label{homological-algebra:lemma-homotopy-invariance}
Homotopic morphisms induce the same maps on cohomology; homotopy equivalences are quasi-isomorphisms. Additive
functors preserve homotopies.
\end{lemma}

\begin{proof}
For a cocycle $z$, $f(z) - g(z) = dh(z) + hd(z) = dh(z)$ is a coboundary. The second statement follows, and the third
because the defining identity is preserved by additive functors.
\end{proof}

\begin{definition}\label{homological-algebra:definition-cone}
The \emph{mapping cone} of $f \colon A^\bullet \to B^\bullet$ is the complex $\mathrm{Cone}(f)$ with
$\mathrm{Cone}(f)^n = A^{n+1} \oplus B^n$ and $d(a, b) = (-da, f(a) + db)$.
\end{definition}

\section{The long exact sequence}\label{homological-algebra:section-long-exact}

\begin{lemma}[Snake lemma]\label{homological-algebra:lemma-snake}
Consider a commutative diagram of modules with exact rows
\begin{center}
\begin{tikzcd}
& A' \arrow[r, "i'"] \arrow[d, "f"] & B' \arrow[r, "p'"] \arrow[d, "g"] & C' \arrow[r] \arrow[d, "h"] & 0 \\
0 \arrow[r] & A \arrow[r, "i"] & B \arrow[r, "p"] & C &
\end{tikzcd}
\end{center}
There is a homomorphism $\delta \colon \ker h \to \operatorname{coker} f$ such that the sequence
\[
\ker f \to \ker g \to \ker h \xrightarrow{\delta} \operatorname{coker} f \to \operatorname{coker} g \to \operatorname{coker} h
\]
is exact. If $i'$ is injective so is $\ker f \to \ker g$, and if $p$ is surjective so is $\operatorname{coker} g \to \operatorname{coker} h$.
\end{lemma}

\begin{proof}
\emph{Definition of $\delta$.} Let $c' \in \ker h$. Choose $b' \in B'$ with $p'(b') = c'$. Then $p(g(b')) = h(c') = 0$, so
$g(b') = i(a)$ for a unique $a \in A$; put $\delta(c') = [a]$. If $b'$ is replaced by $b' + i'(a')$, then $a$ is replaced by
$a + f(a')$, with the same class; so $\delta$ is well defined, and it is clearly a homomorphism.

\emph{Exactness at $\ker g$.} The composite is zero. If $b' \in \ker g$ and $p'(b') = 0$, then $b' = i'(a')$ and
$i(f(a')) = g(b') = 0$, so $f(a') = 0$.

\emph{Exactness at $\ker h$.} If $c' = p'(b')$ with $g(b') = 0$, then $a = 0$ and $\delta(c') = 0$. Conversely if
$\delta(c') = 0$, then $a = f(a'')$, and $b' - i'(a'')$ is a preimage of $c'$ with $g(b' - i'(a'')) = i(a) - i(f(a'')) = 0$.

\emph{Exactness at $\operatorname{coker} f$.} The image of $\delta(c')$ in $\operatorname{coker} g$ is the class of $i(a) = g(b')$, which is
zero. Conversely, if $i(a) = g(b')$ for some $b'$, then $c' = p'(b')$ satisfies $h(c') = p(g(b')) = pi(a) = 0$ and
$\delta(c') = [a]$.

\emph{Exactness at $\operatorname{coker} g$.} The composite is zero. If $b \in B$ maps to zero in $\operatorname{coker} h$, then
$p(b) = h(p'(b'))$ for some $b'$, so $p(b - g(b')) = 0$ and $b - g(b') = i(a)$; hence $[b]$ is the image of $[a]$.

The last two statements are clear.
\end{proof}

\begin{lemma}[Five lemma]\label{homological-algebra:lemma-five}
Consider a commutative diagram of modules with exact rows
$A_1 \to A_2 \to A_3 \to A_4 \to A_5$ and $B_1 \to B_2 \to B_3 \to B_4 \to B_5$ and vertical maps $f_1, \ldots, f_5$. If $f_2$ and
$f_4$ are isomorphisms, $f_1$ is surjective and $f_5$ is injective, then $f_3$ is an isomorphism.
\end{lemma}

\begin{proof}
Write $\alpha_i \colon A_i \to A_{i+1}$ and $\beta_i \colon B_i \to B_{i+1}$. \emph{Injectivity:} let $f_3(x) = 0$. Then
$f_4\alpha_3(x) = 0$, so $\alpha_3(x) = 0$ and $x = \alpha_2(y)$. Now $\beta_2f_2(y) = 0$, so $f_2(y) = \beta_1(z)$ and
$z = f_1(w)$. Then $f_2(y - \alpha_1(w)) = 0$, so $y = \alpha_1(w)$ and $x = \alpha_2\alpha_1(w) = 0$. \emph{Surjectivity:} let
$u \in B_3$. Choose $x \in A_4$ with $f_4(x) = \beta_3(u)$. Then $f_5\alpha_4(x) = \beta_4\beta_3(u) = 0$, so $\alpha_4(x) = 0$ and
$x = \alpha_3(y)$. Then $\beta_3(u - f_3(y)) = 0$, so $u - f_3(y) = \beta_2(v)$ with $v = f_2(t)$, and
$u = f_3(y + \alpha_2(t))$.
\end{proof}

\begin{theorem}[Long exact sequence, or zig-zag lemma]\label{homological-algebra:theorem-long-exact-sequence}
A short exact sequence $0 \to A^\bullet \xrightarrow{f} B^\bullet \xrightarrow{g} C^\bullet \to 0$ of complexes in an abelian category
induces a long exact sequence
\[
\cdots \to H^n(A) \to H^n(B) \to H^n(C) \xrightarrow{\delta} H^{n+1}(A) \to H^{n+1}(B) \to \cdots,
\]
natural in the short exact sequence. For $c \in Z^n(C)$ (in the case of modules), $\delta[c]$ is the class of $a$ where
$f(a) = d(b)$ and $g(b) = c$.
\end{theorem}

\begin{proof}
Applying the snake lemma to the vertical maps $d^{n-1}$ between the short exact sequences in degrees $n - 1$ and $n$, the
cokernel sequence $A^n/B^n(A) \to B^n/B^n(B) \to C^n/B^n(C) \to 0$ is exact; applying it to $d^{n+1}$, the kernel sequence
$0 \to Z^{n+1}(A) \to Z^{n+1}(B) \to Z^{n+1}(C)$ is exact. The differential $d^n$ induces vertical maps from the first of these rows
to the second, whose kernels are the $H^n$ and whose cokernels are the $H^{n+1}$. Apply
Lemma~\ref{homological-algebra:lemma-snake}. Naturality follows from the explicit formula for $\delta$.
\end{proof}

\begin{remark}\label{homological-algebra:remark-zig-zag}
The name \emph{zig-zag lemma} comes from the path that defines the connecting map $\delta$. Start from a cocycle $c \in C^n$; go
up to a preimage $b \in B^n$ with $g(b) = c$; go across by the differential to $d(b) \in B^{n+1}$, which satisfies
$g(d(b)) = d(g(b)) = d(c) = 0$; and go down to the unique $a \in A^{n+1}$ with $f(a) = d(b)$. Then $\delta[c] = [a]$:
\begin{center}
\begin{tikzcd}
& b \arrow[r, mapsto, "g"] \arrow[d, mapsto, "d"'] & c \arrow[d, mapsto, "d"] \\
a \arrow[r, mapsto, "f"] & d(b) \arrow[r, mapsto, "g"] & 0
\end{tikzcd}
\end{center}
The element $a$ depends on the choice of $b$, but its class does not; Exercise~\ref{homological-algebra:exercise-zig-zag} checks
this and the exactness of the long sequence directly, by chasing elements through the diagram, without the snake lemma. In
singular homology, where the short exact sequence is $0 \to C_\bullet(A) \to C_\bullet(X) \to C_\bullet(X, A) \to 0$, the recipe says:
a relative cycle is a chain in $X$ whose boundary lies in $A$, and $\delta$ sends it to that boundary.
\end{remark}

\begin{proposition}\label{homological-algebra:proposition-cone}
For $f \colon A^\bullet \to B^\bullet$ there is a short exact sequence
$0 \to B^\bullet \to \mathrm{Cone}(f) \to A[1]^\bullet \to 0$ whose connecting map $H^n(A[1]) = H^{n+1}(A) \to H^{n+1}(B)$ is
$H^{n+1}(f)$. In particular $f$ is a quasi-isomorphism if and only if $\mathrm{Cone}(f)$ is acyclic.
\end{proposition}

\begin{proof}
The maps $b \mapsto (0, b)$ and $(a, b) \mapsto a$ are morphisms of complexes (with the sign convention of
Notation~\ref{homological-algebra:notation-chain-complexes} for $A[1]$). For a cocycle $a \in A^{n+1}$, the element
$(a, 0)$ lifts it, and $d(a, 0) = (-da, f(a)) = (0, f(a))$; so $\delta[a] = [f(a)]$. The last statement follows from the long
exact sequence.
\end{proof}

\section{Projective and injective objects}\label{homological-algebra:section-projective-injective}

\begin{definition}\label{homological-algebra:definition-projective-injective}
An object $P$ of an abelian category is \emph{projective} if $\Hom(P, -)$ is exact, that is, every morphism $P \to Y$
lifts along every epimorphism $X \to Y$. An object $I$ is \emph{injective} if $\Hom(-, I)$ is exact, that is, every morphism
$X \to I$ extends along every monomorphism $X \to Y$. The category has \emph{enough projectives} (\emph{injectives}) if
every object is a quotient of a projective (a subobject of an injective).
\end{definition}

\begin{proposition}\label{homological-algebra:proposition-projective-modules}
An $R$-module is projective if and only if it is a direct summand of a free module. $R\text{-}\cat{Mod}$ has enough
projectives.
\end{proposition}

\begin{proof}
Free modules are projective by Proposition~\ref{rings:proposition-free-module} (choose lifts of the images of a basis),
and direct summands of projectives are projective. If $P$ is projective, choose a surjection $F \to P$ from a free module;
the identity of $P$ lifts to a section, so $P$ is a direct summand of $F$ by Proposition~\ref{rings:proposition-splitting}.
Every module is a quotient of a free module.
\end{proof}

\begin{proposition}[Baer's criterion]\label{homological-algebra:proposition-baer}
An $R$-module $I$ is injective if and only if every homomorphism $J \to I$ from a left ideal $J \subseteq R$ extends to $R$.
\end{proposition}

\begin{proof}
Necessity is clear. Let $N \subseteq M$ and $\varphi \colon N \to I$. By Zorn's lemma choose a maximal extension
$\psi \colon N' \to I$ with $N \subseteq N' \subseteq M$. If $m \in M \setminus N'$, let $J = \{r : rm \in N'\}$ and extend
$r \mapsto \psi(rm)$ to $\theta \colon R \to I$. Then $n' + rm \mapsto \psi(n') + \theta(r)$ is a well-defined extension of $\psi$ to
$N' + Rm$ (if $n' + rm = 0$ then $r \in J$ and $\theta(r) = \psi(rm) = -\psi(n')$), contradicting maximality.
\end{proof}

\begin{theorem}\label{homological-algebra:theorem-enough-injectives}
$R\text{-}\cat{Mod}$ has enough injectives. An abelian group $D$ is injective if and only if it is \emph{divisible}: for every
$x \in D$ and $n \geq 1$ there is $y$ with $ny = x$.
\end{theorem}

\begin{proof}
By Baer's criterion and Proposition~\ref{groups:proposition-subgroups-integers}, $D$ is injective iff every homomorphism
$n\ZZ \to D$ extends to $\ZZ$, iff divisible. Every abelian group $M$ embeds in a divisible group: write $M = F/K$ with
$F = \ZZ^{(S)}$ free; then $M \subseteq \QQ^{(S)}/K$, which is divisible.

Let $M$ be an $R$-module and $M \subseteq D$ an embedding of abelian groups with $D$ divisible. Make $\Hom_\ZZ(R, D)$ an
$R$-module by $(r\varphi)(s) = \varphi(sr)$. For every $R$-module $N$ there is a natural bijection
$\Hom_R(N, \Hom_\ZZ(R, D)) \cong \Hom_\ZZ(N, D)$, $\psi \mapsto (n \mapsto \psi(n)(1))$, with inverse
$\chi \mapsto (n \mapsto (s \mapsto \chi(sn)))$. The right side is exact in $N$ since $D$ is injective, so $\Hom_\ZZ(R, D)$ is an
injective $R$-module. The map $M \to \Hom_\ZZ(R, D)$, $m \mapsto (s \mapsto sm)$, is $R$-linear and injective.
\end{proof}

\section{Derived functors}\label{homological-algebra:section-derived-functors}

In this section $\mathcal{A}$ and $\mathcal{B}$ are abelian categories.

\begin{definition}\label{homological-algebra:definition-resolution}
An \emph{injective resolution} of an object $X$ is an exact sequence $0 \to X \to I^0 \to I^1 \to \cdots$ with all $I^n$
injective; we regard $I^\bullet$ as a complex with a quasi-isomorphism $X \to I^\bullet$ ($X$ in degree $0$). Projective
resolutions $\cdots \to P_1 \to P_0 \to X \to 0$ are defined dually. If $\mathcal{A}$ has enough injectives, every object has
an injective resolution: embed $X$ in $I^0$, embed $I^0/X$ in $I^1$, and so on.
\end{definition}

\begin{theorem}[Comparison theorem]\label{homological-algebra:theorem-comparison}
Let $0 \to X \to E^\bullet$ be exact with $E^n = 0$ for $n < 0$, and let $0 \to Y \to I^\bullet$ be a complex with all $I^n$ injective,
$I^n = 0$ for $n < 0$. Every morphism $\varphi \colon X \to Y$ extends to a morphism of complexes $E^\bullet \to I^\bullet$, and any two
extensions are homotopic.
\end{theorem}

\begin{proof}
Write $\varepsilon \colon X \to E^0$ and $\eta \colon Y \to I^0$. The morphism $\eta\varphi$ extends along the monomorphism $\varepsilon$ to
$f^0 \colon E^0 \to I^0$. Suppose $f^0, \ldots, f^n$ are constructed commuting with the differentials (and with $\varepsilon$, $\eta$).
Then $d f^n$ vanishes on the image of $E^{n-1} \to E^n$ (or of $X$ when $n = 0$), so it factors through
$E^n/\im(E^{n-1}) \hookrightarrow E^{n+1}$, and extends along this monomorphism to $f^{n+1} \colon E^{n+1} \to I^{n+1}$ by injectivity.

If $f$, $g$ both extend $\varphi$, let $u = f - g$, which extends $0$, so $u^0\varepsilon = 0$. Put $h^0 = 0 \colon E^0 \to I^{-1} = 0$.
Suppose $h^0, \ldots, h^n$ are constructed with $u^k = dh^k + h^{k+1}d$ for all $k < n$. For $n \geq 1$,
\[
(u^n - dh^n)d = u^nd - dh^nd = du^{n-1} - d(u^{n-1} - dh^{n-1}) = ddh^{n-1} = 0,
\]
using that $u$ commutes with the differentials and the identity for $k = n - 1$; for $n = 0$, $(u^0 - dh^0)\varepsilon = u^0\varepsilon = 0$.
So $u^n - dh^n$ vanishes on the image of $E^{n-1} \to E^n$ (of $X \to E^0$ if $n = 0$), factors through the monomorphism
$E^n/\im \hookrightarrow E^{n+1}$, and extends to $h^{n+1} \colon E^{n+1} \to I^n$ with $h^{n+1}d = u^n - dh^n$.
\end{proof}

\begin{lemma}[Horseshoe lemma]\label{homological-algebra:lemma-horseshoe}
Let $0 \to X' \to X \to X'' \to 0$ be exact, and let $X' \to I'^\bullet$, $X'' \to I''^\bullet$ be injective resolutions. There is an
injective resolution $X \to I^\bullet$ with $I^n = I'^n \oplus I''^n$ such that $0 \to I'^\bullet \to I^\bullet \to I''^\bullet \to 0$ (with the
inclusion and projection in each degree) is a short exact sequence of complexes compatible with the given sequence.
\end{lemma}

\begin{proof}
Let $\varepsilon' \colon X' \to I'^0$, $\varepsilon'' \colon X'' \to I''^0$, and $\iota \colon X' \to X$, $\pi \colon X \to X''$. Extend $\varepsilon'$
along $\iota$ to $\lambda \colon X \to I'^0$ and define $\varepsilon = (\lambda, \varepsilon''\pi) \colon X \to I'^0 \oplus I''^0$. The diagram
$0 \to X' \to X \to X'' \to 0$ over $0 \to I'^0 \to I^0 \to I''^0 \to 0$ commutes, and the snake lemma shows that $\varepsilon$ is a
monomorphism and that $0 \to \operatorname{coker}\varepsilon' \to \operatorname{coker}\varepsilon \to \operatorname{coker}\varepsilon'' \to 0$ is exact. Apply the same
construction to this sequence and the resolutions $I'^{\geq 1}$, $I''^{\geq 1}$ of the cokernels, and continue by induction.
\end{proof}

\begin{definition}\label{homological-algebra:definition-derived-functor}
Let $\mathcal{A}$ have enough injectives and $F \colon \mathcal{A} \to \mathcal{B}$ be a left exact functor. For each $X$ choose an
injective resolution $X \to I_X^\bullet$ and define the \emph{right derived functors} $R^nF(X) = H^n(F(I_X^\bullet))$. A morphism
$X \to Y$ induces a morphism of resolutions, unique up to homotopy (Theorem~\ref{homological-algebra:theorem-comparison}),
hence a well-defined morphism $R^nF(X) \to R^nF(Y)$ by Lemma~\ref{homological-algebra:lemma-homotopy-invariance}. Left derived
functors $L_nF$ of a right exact functor on a category with enough projectives are defined dually using projective
resolutions.
\end{definition}

\begin{theorem}\label{homological-algebra:theorem-derived-functors}
\begin{enumerate}
\item The $R^nF$ are additive functors, independent of the choice of resolutions up to canonical isomorphism, with
$R^nF = 0$ for $n < 0$ and $R^0F \cong F$.
\item $R^nF(I) = 0$ for $n > 0$ and $I$ injective.
\item Every short exact sequence $0 \to X' \to X \to X'' \to 0$ induces a long exact sequence
$0 \to F(X') \to F(X) \to F(X'') \to R^1F(X') \to R^1F(X) \to \cdots$, natural in the short exact sequence.
\end{enumerate}
\end{theorem}

\begin{proof}
(1) Two resolutions of $X$ are related by morphisms extending $\id_X$ in both directions, whose composites are homotopic
to the identities by the uniqueness part of Theorem~\ref{homological-algebra:theorem-comparison}; applying $F$ gives
homotopy equivalences and canonical isomorphisms on cohomology. Functoriality and additivity follow from uniqueness up
to homotopy. $R^0F(X) = \ker(F(I^0) \to F(I^1)) = F(\ker(I^0 \to I^1)) = F(X)$ by left exactness.
(2) Use the resolution $0 \to I \to I \to 0$.
(3) Choose resolutions of $X'$ and $X''$ and form the resolution of $X$ given by
Lemma~\ref{homological-algebra:lemma-horseshoe}. The sequence $0 \to I'^\bullet \to I^\bullet \to I''^\bullet \to 0$ is split in every
degree, so it stays exact after applying the additive functor $F$. Apply
Theorem~\ref{homological-algebra:theorem-long-exact-sequence}. Naturality follows from the comparison theorem.
\end{proof}

\begin{proposition}\label{homological-algebra:proposition-acyclic-resolutions}
Call $J$ \emph{$F$-acyclic} if $R^nF(J) = 0$ for all $n > 0$. If $0 \to X \to J^0 \to J^1 \to \cdots$ is exact with all $J^n$
$F$-acyclic, then $R^nF(X) \cong H^n(F(J^\bullet))$ for all $n$.
\end{proposition}

\begin{proof}
Let $K = \operatorname{coker}(X \to J^0) = \ker(J^1 \to J^2)$, so that $0 \to X \to J^0 \to K \to 0$ is exact and $K \to J^{\geq 1}$ is an $F$-acyclic
resolution. The long exact sequence and acyclicity of $J^0$ give $R^nF(X) \cong R^{n-1}F(K)$ for $n \geq 2$, and
$R^1F(X) \cong \operatorname{coker}(F(J^0) \to F(K))$. By left exactness $F(K) = \ker(F(J^1) \to F(J^2))$, so
$R^1F(X) \cong H^1(F(J^\bullet))$. For $n \geq 2$ argue by induction on $n$, for all objects and resolutions simultaneously:
$K \to J^1 \to J^2 \to \cdots$ is an $F$-acyclic resolution of $K$, with $J^1$ in degree $0$, so
$R^nF(X) \cong R^{n-1}F(K) \cong H^{n-1}(F(J^1) \to F(J^2) \to \cdots) = H^n(F(J^\bullet))$. The case $n = 0$ is left exactness.
\end{proof}

\section{Double complexes}\label{homological-algebra:section-double-complexes}

\begin{definition}\label{homological-algebra:definition-double-complex}
A \emph{double complex} $C^{\bullet,\bullet}$ consists of objects $C^{p,q}$ with morphisms
$d_h \colon C^{p,q} \to C^{p+1,q}$ and $d_v \colon C^{p,q} \to C^{p,q+1}$ such that $d_h^2 = 0$, $d_v^2 = 0$ and $d_hd_v = d_vd_h$. It is
\emph{first quadrant} if $C^{p,q} = 0$ unless $p, q \geq 0$. Its \emph{total complex} is
$\mathrm{Tot}^n = \bigoplus_{p+q=n}C^{p,q}$ with $d = d_h + (-1)^pd_v$ on $C^{p,q}$; then $d^2 = 0$.
\end{definition}

\begin{lemma}\label{homological-algebra:lemma-acyclic-assembly}
Let $C$ be a first-quadrant double complex of modules whose rows are exact in positive degrees: for each $q$, the complex
$C^{0,q} \to C^{1,q} \to \cdots$ is exact at $C^{p,q}$ for all $p \geq 1$. Let $K^q = \ker(d_h \colon C^{0,q} \to C^{1,q})$, a complex with
differential $d_v$. Then the inclusion $K^\bullet \to \mathrm{Tot}^\bullet(C)$ is a quasi-isomorphism.
\end{lemma}

\begin{proof}
The inclusion is a morphism of complexes, since $d = d_v$ on $K$. \emph{Surjectivity on cohomology.} Let
$x = (x^{0,n}, x^{1,n-1}, \ldots, x^{n,0})$ be a cocycle of degree $n$. Let $p$ be largest with $x^{p,n-p} \neq 0$. If $p \geq 1$, the
component of $dx$ in $C^{p+1,n-p}$ is $d_hx^{p,n-p} = 0$ (the other contribution comes from $x^{p+1, n-p-1} = 0$), so
$x^{p,n-p} = d_hy$ with $y \in C^{p-1,n-p}$, and $x - dy$ is a cohomologous cocycle whose components vanish in positions
$\geq p$. Repeating, we reach a cocycle concentrated in $C^{0,n}$; the vanishing of its $C^{1,n}$-component of $d$ means it
lies in $K^n$, and it is a $d_v$-cocycle.

\emph{Injectivity on cohomology.} Let $k \in K^n$ with $k = dy$ for $y \in \mathrm{Tot}^{n-1}$. As above, if $p \geq 1$ is largest with
$y^{p,n-1-p} \neq 0$, the component of $dy$ in $C^{p+1,n-1-p}$ is $d_hy^{p,n-1-p}$, which is zero because $dy = k$ lives in
$C^{0,n}$ (here $p + 1 \geq 2 > 0$); so $y^{p,n-1-p} = d_hz$ and we replace $y$ by $y - dz$, which does not change $dy$. We reach
$y \in C^{0,n-1}$ with $d_hy$ equal to the $C^{1,n-1}$-component of $k$, which is $0$; so $y \in K^{n-1}$ and $k = d_vy$.
\end{proof}

\section{Ext and Tor}\label{homological-algebra:section-ext-tor}

In this section $R$ is a ring and modules are left $R$-modules; for $\Tor$ we assume $R$ commutative.

\begin{definition}\label{homological-algebra:definition-ext-tor}
$\Ext^n_R(M, N) = R^n\Hom_R(M, -)(N)$, computed with an injective resolution of $N$, and
$\Tor_n^R(M, N) = L_n(M \otimes_R -)(N)$, computed with a projective resolution of $N$.
\end{definition}

\begin{theorem}[Balance]\label{homological-algebra:theorem-balance}
If $P_\bullet \to M$ is a projective resolution, then $\Ext^n_R(M, N) \cong H^n(\Hom_R(P_\bullet, N))$. Similarly
$\Tor_n^R(M, N) \cong H_n(P_\bullet \otimes_R N)$, and $\Tor_n^R(M, N) \cong \Tor_n^R(N, M)$. Consequently $\Ext^n(-, N)$ are the right
derived functors of the contravariant left exact functor $\Hom_R(-, N)$, and there are long exact sequences in both
variables.
\end{theorem}

\begin{proof}
Let $N \to I^\bullet$ be an injective resolution and $C^{p,q} = \Hom_R(P_p, I^q)$, a first-quadrant double complex. For fixed $q$ the
row $\Hom_R(P_\bullet, I^q)$ is exact in positive degrees with kernel $\Hom_R(M, I^q)$ in degree $0$, because $\Hom_R(-, I^q)$ is
exact and $P_\bullet \to M \to 0$ is exact. For fixed $p$ the column $\Hom_R(P_p, I^\bullet)$ is exact in positive degrees with kernel
$\Hom_R(P_p, N)$, because $P_p$ is projective. By Lemma~\ref{homological-algebra:lemma-acyclic-assembly}, applied to $C$ and to its
transpose, both $\Hom_R(M, I^\bullet)$ and $\Hom_R(P_\bullet, N)$ are quasi-isomorphic to $\mathrm{Tot}(C)$. The same argument with
$P_p \otimes Q_q$ for projective resolutions $P$ of $M$ and $Q$ of $N$ (projective modules are flat, being summands of free
modules) gives the statements for $\Tor$, and symmetry follows from $P_p \otimes Q_q \cong Q_q \otimes P_p$. The long exact sequence
in the first variable follows from the horseshoe lemma for projective resolutions and
Theorem~\ref{homological-algebra:theorem-long-exact-sequence}.
\end{proof}

\begin{example}\label{homological-algebra:example-ext-tor-cyclic}
Using the projective resolution $0 \to \ZZ \xrightarrow{n} \ZZ \to \ZZ/n\ZZ \to 0$, for every abelian group $A$:
$\Hom(\ZZ/n\ZZ, A) = A[n] = \{a : na = 0\}$, $\Ext^1_\ZZ(\ZZ/n\ZZ, A) = A/nA$, $\ZZ/n\ZZ \otimes A = A/nA$ and
$\Tor_1^\ZZ(\ZZ/n\ZZ, A) = A[n]$; all higher groups vanish. For instance $\Ext^1_\ZZ(\ZZ/n\ZZ, \ZZ) \cong \ZZ/n\ZZ$, although
$\Hom(\ZZ/n\ZZ, \ZZ) = 0$.
\end{example}

\begin{proposition}\label{homological-algebra:proposition-flat-tor}
A module $N$ is flat if and only if $\Tor_1^R(M, N) = 0$ for all $M$; then $\Tor_n^R(M, N) = 0$ for all $n \geq 1$. If $R$ is
Noetherian local with residue field $k$ and $N$ is finitely generated, then $N$ is flat if and only if it is free, if and only if
$\Tor_1^R(k, N) = 0$.
\end{proposition}

\begin{proof}
If $0 \to M' \to M \to M'' \to 0$ is exact, the long exact sequence
$\Tor_1(M'', N) \to M' \otimes N \to M \otimes N$ shows that vanishing of $\Tor_1$ gives injectivity; conversely if $N$ is flat,
computing $\Tor$ with a projective resolution of $M$ tensored with $N$ (Theorem~\ref{homological-algebra:theorem-balance})
gives an exact complex. For the second statement, free implies flat. Suppose $\Tor_1(k, N) = 0$. Choose a minimal set of
generators, giving $0 \to K \to R^n \to N \to 0$ with $n = \dim_k N/\mathfrak{m}N$
(Corollary~\ref{commutative-algebra:corollary-nakayama-generators}). Tensoring with $k$ gives the exact sequence
$0 \to K \otimes k \to k^n \to N \otimes k \to 0$, in which the last map is an isomorphism of $n$-dimensional spaces; so
$K/\mathfrak{m}K = 0$, and $K = 0$ by Nakayama's lemma, as $K$ is finitely generated ($R$ is Noetherian).
\end{proof}

\begin{proposition}\label{homological-algebra:proposition-pid-submodules-free}
If $R$ is a principal ideal domain, every submodule of a free $R$-module is free. Hence every $R$-module has a projective
resolution of length at most $1$, and $\Ext^n_R = \Tor_n^R = 0$ for $n \geq 2$.
\end{proposition}

\begin{proof}
Let $F$ be free with basis $(e_i)_{i \in I}$, where $I$ is well-ordered (Theorem~\ref{sets:theorem-choice-equivalents}), and
$M \subseteq F$. For $i \in I$ let $F_{\leq i}$ and $F_{<i}$ be spanned by the $e_j$ with $j \leq i$, respectively $j < i$, let
$M_i = M \cap F_{\leq i}$, and let $\pi_i \colon M_i \to R$ be the $e_i$-coordinate. Its image is an ideal $(a_i)$; if $a_i \neq 0$ choose
$m_i \in M_i$ with $\pi_i(m_i) = a_i$. The $m_i$ are linearly independent: in a relation $\sum c_im_i = 0$, the largest $i$ with
$c_i \neq 0$ gives $c_ia_i = 0$ on the $e_i$-coordinate. They generate $M$: otherwise let $i$ be least such that $M_i$ contains an
element $m$ not in their span. Then $\pi_i(m) = ra_i$, and $m - rm_i$ (or $m$ itself if $a_i = 0$) lies in $M \cap F_{<i}$; having
finite support, it lies in some $M_j$ with $j < i$, hence in the span, and so does $m$. For the last statement, write $M$ as a
quotient of a free module; the kernel is free.
\end{proof}

\begin{proposition}\label{homological-algebra:proposition-ext-extensions}
Let $M$, $N$ be $R$-modules. There is a bijection between $\Ext^1_R(M, N)$ and the set of equivalence classes of short exact
sequences $0 \to N \to E \to M \to 0$, where two sequences are equivalent if there is an isomorphism of the middle terms compatible
with the given maps; the zero element corresponds to the split sequence.
\end{proposition}

\begin{proof}
Choose $0 \to K \xrightarrow{j} P \to M \to 0$ with $P$ projective. By Theorem~\ref{homological-algebra:theorem-balance} and the long
exact sequence, $\Ext^1(M, N) = \operatorname{coker}(\Hom(P, N) \to \Hom(K, N))$. Given an extension $0 \to N \to E \xrightarrow{q} M \to 0$, lift
$P \to M$ to $\beta \colon P \to E$; then $\beta j$ lands in $N$, and its class $\theta(E)$ does not depend on $\beta$ (two lifts differ by a
map $P \to N$). Conversely, given $\alpha \colon K \to N$, let $E_\alpha$ be the pushout of $N \xleftarrow{\alpha} K \xrightarrow{j} P$, that is,
$(N \oplus P)/\{(\alpha(k), -j(k)) : k \in K\}$; it sits in an exact sequence $0 \to N \to E_\alpha \to M \to 0$, and $\alpha$ and
$\alpha + \gamma j$ ($\gamma \colon P \to N$) give equivalent extensions via $E_\alpha \to E_{\alpha + \gamma j}$, $(n, p) \mapsto (n - \gamma(p), p)$, which maps
the relation $(\alpha(k), -j(k))$ to $((\alpha + \gamma j)(k), -j(k))$. The lift $p \mapsto (0, p)$ of $P \to M$ to $E_\alpha$ restricts on $K$ to
$k \mapsto (0, j(k)) = (\alpha(k), 0)$, so $\theta(E_\alpha)$ is the class of $\alpha$. For an extension $E$ with lift $\beta$, the map
$E_{\beta j} \to E$, $(n, p) \mapsto n + \beta(p)$, is well defined and an equivalence by the five lemma. Hence the constructions are
inverse. The split extension admits a lift with $\beta j = 0$.
\end{proof}

\begin{example}\label{homological-algebra:example-ext-zp}
Let $p$ be a prime. By Proposition~\ref{homological-algebra:proposition-ext-extensions}, computed with the resolution $0 \to p\ZZ \xrightarrow{j} \ZZ \to \ZZ/p \to 0$: every $f \colon \ZZ \to \ZZ/p$ vanishes on $p\ZZ$, so
$\Ext^1_\ZZ(\ZZ/p, \ZZ/p) = \Hom(p\ZZ, \ZZ/p) \cong \ZZ/p$, via $\alpha \mapsto a = \alpha(p)$. For $a = 0$ the extension is split, $E = \ZZ/p \oplus \ZZ/p$. For $a \neq 0$ the pushout is
$E_\alpha = (\ZZ/p \oplus \ZZ)/\langle(a, -p)\rangle$, in which $(0, 1)$ has order $p^2$: if $m(0, 1) = k(a, -p)$, then $ka \equiv 0 \pmod p$ forces $p \mid k$ and $p^2 \mid m$. So $E_\alpha \cong \ZZ/p^2$, with $\ZZ/p$
embedded by $n \mapsto na^{-1}p$. The $p - 1$ nonzero classes thus give pairwise inequivalent extensions whose middle terms are all isomorphic to $\ZZ/p^2$: equivalence of extensions remembers the maps, not
only the middle group.
\end{example}

\begin{counterexample}\label{homological-algebra:counterexample-flat-projective-injective}
Projective, flat and injective are different conditions, already for abelian groups.
\begin{enumerate}
\item $\QQ$ is flat, since $\QQ = S^{-1}\ZZ$ with $S = \ZZ \setminus \{0\}$ (Proposition~\ref{commutative-algebra:proposition-localisation}(2)), and injective, since it is divisible
(Theorem~\ref{homological-algebra:theorem-enough-injectives}). It is not projective: a projective abelian group is a direct summand, hence a subgroup, of a free group, hence free
(Propositions~\ref{homological-algebra:proposition-projective-modules} and~\ref{homological-algebra:proposition-pid-submodules-free}); but $\Hom(\QQ, \ZZ) = 0$, because the image of $1$ under a homomorphism
$\QQ \to \ZZ$ is divisible by every $n \geq 1$ in $\ZZ$, while a nonzero free group has a nonzero coordinate map to $\ZZ$.
\item $\ZZ$ is projective, hence flat, but not injective, as it is not divisible.
\item $\ZZ/n$ for $n \geq 2$ is none of the three: it is not flat, since $\Tor_1(\ZZ/n, \ZZ/n) = \ZZ/n \neq 0$ (Example~\ref{homological-algebra:example-ext-tor-cyclic} and
Proposition~\ref{homological-algebra:proposition-flat-tor}), hence not projective, and it is not divisible, hence not injective.
\end{enumerate}
\end{counterexample}

\section{K\"unneth and universal coefficient theorems}\label{homological-algebra:section-kunneth}

In this section $R$ is a principal ideal domain and we use chain complexes of $R$-modules. The tensor product of chain
complexes is $(C \otimes D)_n = \bigoplus_{p+q=n}C_p \otimes D_q$ with $d(x \otimes y) = dx \otimes y + (-1)^px \otimes dy$ for $x \in C_p$, and
$\Hom(C, M)$ is the cochain complex $\Hom(C_n, M)$ with $(d\varphi)(x) = (-1)^{n+1}\varphi(dx)$ for $\varphi \in \Hom(C_n, M)$.

\begin{theorem}[K\"unneth formula]\label{homological-algebra:theorem-kunneth}
Let $C$ be a chain complex of free $R$-modules and $D$ any chain complex. For every $n$ there is a natural short exact sequence
\[
0 \to \bigoplus_{p+q=n}H_p(C) \otimes H_q(D) \to H_n(C \otimes D) \to \bigoplus_{p+q=n-1}\Tor_1^R(H_p(C), H_q(D)) \to 0,
\]
where the first map sends $[x] \otimes [y]$ to $[x \otimes y]$. If $D$ is also a complex of free modules, the sequence splits
(non-naturally).
\end{theorem}

\begin{proof}
Let $Z_p = \ker d_p$ and $B_p = \im d_{p+1} \subseteq C_p$; they are free by
Proposition~\ref{homological-algebra:proposition-pid-submodules-free}, hence flat. Regard $Z$ and $B$ as complexes with zero
differential, and let $B'$ be the complex with $B'_p = B_{p-1}$ and zero differential. The sequences
$0 \to Z_p \to C_p \xrightarrow{d} B_{p-1} \to 0$ form a short exact sequence of complexes $0 \to Z \to C \to B' \to 0$, split in each degree
since $B_{p-1}$ is free; so it stays exact after tensoring with $D$. As $Z$ and $B$ are
flat with zero differential, $H_n(Z \otimes D) = \bigoplus Z_p \otimes H_q(D)$ and similarly for $B$. The long exact sequence reads
\[
\cdots \to \bigoplus_{p+q=n}B_p \otimes H_q(D) \xrightarrow{\partial} \bigoplus_{p+q=n}Z_p \otimes H_q(D) \to H_n(C \otimes D) \to
\bigoplus_{p+q=n-1}B_p \otimes H_q(D) \xrightarrow{\partial} \cdots,
\]
and computing the connecting map with Theorem~\ref{homological-algebra:theorem-long-exact-sequence} shows that $\partial$ is induced by
the inclusions $B_p \to Z_p$ (a class $b \otimes [y]$ lifts to $c \otimes y$ with $dc = b$, and $d(c \otimes y) = b \otimes y$). Hence
$H_n(C \otimes D)$ is an extension of $\ker(\partial)$ in degree $n - 1$ by $\operatorname{coker}(\partial)$ in degree $n$. Since
$0 \to B_p \to Z_p \to H_p(C) \to 0$ is a free resolution, $\operatorname{coker}(B_p \otimes H_q(D) \to Z_p \otimes H_q(D)) = H_p(C) \otimes H_q(D)$ and
$\ker = \Tor_1(H_p(C), H_q(D))$, by right exactness of $\otimes$ and Theorem~\ref{homological-algebra:theorem-balance}.

If $D$ is free as well, choose retractions $r_C \colon C_p \to Z_p$ (which exist since $B_{p-1}$ is free) and $r_D$. The composites
$C_p \to Z_p \to H_p(C)$ form a morphism of complexes $C \to H(C)$ (to the complex with zero differential), since $r_C(dx) = dx \in B_p$
maps to $0$; similarly $D \to H(D)$. Their tensor product induces $H_n(C \otimes D) \to \bigoplus H_p(C) \otimes H_q(D)$, which is a left
inverse of the first map.
\end{proof}

\begin{theorem}[Universal coefficient theorem]\label{homological-algebra:theorem-universal-coefficients}
Let $C$ be a chain complex of free $R$-modules and $M$ an $R$-module. There is a natural short exact sequence
\[
0 \to \Ext^1_R(H_{n-1}(C), M) \to H^n(\Hom(C, M)) \to \Hom_R(H_n(C), M) \to 0,
\]
where the second map sends the class of a cocycle $\varphi$ to its restriction to cycles. It splits non-naturally.
\end{theorem}

\begin{proof}
Apply $\Hom(-, M)$ to the degreewise split sequence $0 \to Z \to C \to B' \to 0$ of the previous proof. As $Z$ and $B$ have zero
differential, the long exact sequence becomes
$\cdots \to \Hom(B_{n-1}, M) \to H^n(\Hom(C, M)) \to \Hom(Z_n, M) \xrightarrow{\partial} \Hom(B_n, M) \to \cdots$, where $\partial$ is restriction
along $B_n \subseteq Z_n$ (up to sign). From the free resolution $0 \to B_n \to Z_n \to H_n(C) \to 0$ we get
$\ker(\partial) = \Hom(H_n(C), M)$ and $\operatorname{coker}(\Hom(Z_{n-1}, M) \to \Hom(B_{n-1}, M)) = \Ext^1(H_{n-1}(C), M)$, which gives the
sequence. A splitting is obtained from a retraction $r \colon C_n \to Z_n$: a homomorphism $H_n(C) \to M$ composed with
$C_n \to Z_n \to H_n(C)$ is a cocycle.
\end{proof}

\begin{example}\label{homological-algebra:example-uct}
Let $C$ be the complex $\cdots \to 0 \to \ZZ \xrightarrow{2} \ZZ \to 0$ in degrees $1$ and $0$, so $H_0 = \ZZ/2\ZZ$ and $H_1 = 0$. Then
$\Hom(C, \ZZ)$ is $\ZZ \xrightarrow{\pm 2} \ZZ$ in degrees $0$ and $1$, with $H^0 = 0$ and $H^1 = \ZZ/2\ZZ$. Indeed
$\Hom(H_1, \ZZ) = 0$ and $\Ext^1(H_0, \ZZ) = \ZZ/2\ZZ$: torsion in homology shifts up one degree in cohomology. This is the
algebraic reason why $H^2(\RR\PP^2, \ZZ) \cong \ZZ/2\ZZ$ while $H_2(\RR\PP^2, \ZZ) = 0$ (Chapter~\ref{singular-homology}).
\end{example}

\section{Spectral Sequences}\label{homological-algebra:section-spectral-sequences}
A spectral sequence is a bookkeeping device for computing the cohomology of a filtered complex from the
cohomology of its graded pieces, in a sequence of successive approximations. This chapter constructs the
spectral sequence of a filtered complex with complete proofs, derives the two spectral sequences of a double
complex, and proves the Grothendieck spectral sequence of a composition of functors. Later chapters use the
Leray spectral sequence (Chapter~\ref{sheaves}), the Fr\"olicher spectral sequence
(Chapter~\ref{complex-manifolds}) and the weight and Hodge spectral sequences of mixed Hodge theory
(Chapter~\ref{mixed-hodge-structures}).

References are \cite[Chapter 5]{weibel-homological} and \cite{mccleary-users-guide}. As in
Chapter~\ref{homological-algebra}, we argue with elements; by
Remark~\ref{homological-algebra:remark-freyd-mitchell} the results hold in every abelian category.

\noindent\textbf{Prerequisites.} Chapter~\ref{homological-algebra} (Homological Algebra).

\section{Spectral sequences}\label{homological-algebra:section-definitions}

\begin{definition}\label{homological-algebra:definition-spectral-sequence}
A (cohomological) \emph{spectral sequence} starting at page $r_0$ consists of objects $E_r^{p,q}$ for $p, q \in \ZZ$
and $r \geq r_0$, morphisms $d_r \colon E_r^{p,q} \to E_r^{p+r, q-r+1}$ with $d_r \circ d_r = 0$, and isomorphisms
$E_{r+1}^{p,q} \cong \ker(d_r)/\im(d_r)$ at $E_r^{p,q}$. The object $E_r = \bigoplus_{p,q}E_r^{p,q}$ is the $r$-th \emph{page}, and
$p + q$ is the \emph{total degree}. It is a \emph{first-quadrant} spectral sequence if $E_{r_0}^{p,q} = 0$ unless $p, q \geq 0$.
\end{definition}

\begin{lemma}\label{homological-algebra:lemma-first-quadrant-stable}
In a first-quadrant spectral sequence $E_r^{p,q} = 0$ unless $p, q \geq 0$, for all $r$, and $E_r^{p,q} = E_{r+1}^{p,q}$ whenever
$r > \max(p, q + 1)$. The common value is denoted $E_\infty^{p,q}$.
\end{lemma}

\begin{proof}
The first claim holds because $E_{r+1}^{p,q}$ is a subquotient of $E_r^{p,q}$. For $r > p$ the source $E_r^{p-r, q+r-1}$ of the incoming
differential has negative first index, and for $r > q + 1$ the target $E_r^{p+r, q-r+1}$ of the outgoing differential has negative
second index; so both differentials vanish.
\end{proof}

\begin{definition}\label{homological-algebra:definition-convergence}
A first-quadrant spectral sequence \emph{converges} to objects $H^n$, written $E_{r_0}^{p,q} \Rightarrow H^{p+q}$, if each $H^n$ has a
finite decreasing filtration $H^n = F^0H^n \supseteq F^1H^n \supseteq \cdots \supseteq F^{n+1}H^n = 0$ with isomorphisms
$E_\infty^{p,q} \cong F^pH^{p+q}/F^{p+1}H^{p+q}$. The spectral sequence \emph{degenerates} at page $r$ if $d_s = 0$ for all
$s \geq r$, so that $E_r = E_\infty$.
\end{definition}

\begin{proposition}\label{homological-algebra:proposition-edge-maps}
Let $E_2^{p,q} \Rightarrow H^{p+q}$ be a convergent first-quadrant spectral sequence.
\begin{enumerate}
\item There are \emph{edge maps} $E_2^{n,0} \to H^n$ and $H^n \to E_2^{0,n}$: the first is the composite
$E_2^{n,0} \twoheadrightarrow E_\infty^{n,0} = F^nH^n \hookrightarrow H^n$, the second $H^n \twoheadrightarrow H^n/F^1H^n = E_\infty^{0,n} \hookrightarrow E_2^{0,n}$.
\item (Five-term exact sequence) The sequence
$0 \to E_2^{1,0} \to H^1 \to E_2^{0,1} \xrightarrow{d_2} E_2^{2,0} \to H^2$ is exact.
\item If $E_2^{p,q} = 0$ for $q \neq 0$, then $H^n \cong E_2^{n,0}$. If $E_2^{p,q} = 0$ for $p \neq 0$, then $H^n \cong E_2^{0,n}$.
\end{enumerate}
\end{proposition}

\begin{proof}
(1) Since nothing is in negative degrees, $E_{r+1}^{n,0}$ is a quotient of $E_r^{n,0}$ (no outgoing differential lands in $q < 0$) and
$E_{r+1}^{0,n}$ a subobject of $E_r^{0,n}$ (no incoming differential starts at $p < 0$). Also $F^{n+1}H^n = 0$, so
$E_\infty^{n,0} = F^nH^n$, and $F^0H^n = H^n$, so $E_\infty^{0,n} = H^n/F^1H^n$.
(2) In total degree $1$, $H^1$ has the filtration $0 \subseteq F^1H^1 \subseteq H^1$ with $F^1H^1 = E_\infty^{1,0} = E_2^{1,0}$ (no differential
reaches or leaves $E_r^{1,0}$ nontrivially: incoming would start at $p = 1 - r < 0$, outgoing lands at $q = 1 - r < 0$) and
$H^1/F^1H^1 = E_\infty^{0,1} = \ker(d_2 \colon E_2^{0,1} \to E_2^{2,0})$ (for $r \geq 3$ all differentials at $(0,1)$ vanish). This gives
exactness at $E_2^{1,0}$, $H^1$ and $E_2^{0,1}$. Next $E_\infty^{2,0} = E_3^{2,0} = E_2^{2,0}/\im(d_2)$, and it injects into $H^2$ as
$F^2H^2$; so the kernel of $E_2^{2,0} \to H^2$ is $\im(d_2)$.
(3) All but one of the graded pieces $E_\infty^{p, n-p}$ of $H^n$ vanish, and all differentials vanish because either the source or the
target is zero.
\end{proof}

\begin{proposition}\label{homological-algebra:proposition-two-rows}
Let $E_2^{p,q} \Rightarrow H^{p+q}$ be a convergent first-quadrant spectral sequence with $E_2^{p,q} = 0$ unless $q \in \{0, s\}$, for some
$s \geq 1$. Then there is a long exact sequence
\[
\cdots \to E_2^{n,0} \to H^n \to E_2^{n-s,s} \xrightarrow{d_{s+1}} E_2^{n+1,0} \to H^{n+1} \to \cdots.
\]
\end{proposition}

\begin{proof}
The only possibly nonzero differentials are $d_{s+1} \colon E_{s+1}^{p,s} \to E_{s+1}^{p+s+1,0}$, and $E_{s+1} = E_2$. So
$E_\infty^{n,0} = \operatorname{coker}(d_{s+1} \colon E_2^{n-s-1,s} \to E_2^{n,0})$ and $E_\infty^{n-s,s} = \ker(d_{s+1} \colon E_2^{n-s,s} \to E_2^{n+1,0})$, and $H^n$ is
an extension $0 \to E_\infty^{n,0} \to H^n \to E_\infty^{n-s,s} \to 0$. Splicing these sequences gives the claim.
\end{proof}

\begin{proposition}\label{homological-algebra:proposition-euler-characteristic}
Let $E$ be a spectral sequence of finite-dimensional vector spaces over a field, with only finitely many nonzero $E_{r_0}^{p,q}$. Then $\chi(E_r) = \sum_{p,q}(-1)^{p+q}\dim E_r^{p,q}$ does not depend on $r$, and if
$E_{r_0}^{p,q} \Rightarrow H^{p+q}$ converges, then $\chi(E_{r_0}) = \sum_n(-1)^n\dim H^n$. In particular $\dim H^n \leq \sum_{p+q=n}\dim E_r^{p,q}$ for every $r$.
\end{proposition}

\begin{proof}
Grading $E_r$ by total degree $n = p + q$, the differential $d_r$ raises $n$ by one, so $(E_r, d_r)$ is a finite complex of finite-dimensional spaces with cohomology $E_{r+1}$. For such a complex
$C^\bullet$, rank--nullity gives $\dim C^n = \dim Z^n + \dim B^{n+1}$ and $\dim H^n = \dim Z^n - \dim B^n$, so $\sum_n(-1)^n\dim C^n = \sum_n(-1)^n\dim H^n$. Hence $\chi(E_{r+1}) = \chi(E_r)$. Only finitely many
positions are nonzero, so the pages stabilise, and convergence gives $\dim H^n = \sum_p\dim E_\infty^{p,n-p}$. The inequality holds because each $E_\infty^{p,q}$ is a subquotient of $E_r^{p,q}$.
\end{proof}

\begin{counterexample}\label{homological-algebra:counterexample-extension}
The limit of a spectral sequence determines the abutment only up to extensions. Let $K$ and $K'$ be the complexes concentrated in degree $0$, with zero differential, given by $K^0 = \ZZ/4$ with the filtration
$F^0 = \ZZ/4 \supseteq F^1 = 2\ZZ/4 \supseteq F^2 = 0$, and $K'^0 = \ZZ/2 \oplus \ZZ/2$ with $F^0 = K'^0 \supseteq F^1 = 0 \oplus \ZZ/2 \supseteq F^2 = 0$. In both cases all differentials vanish and the spectral
sequence of the filtered complex (Section~\ref{homological-algebra:section-filtered-complex}) has $E_0^{0,0} = \ZZ/2$ and $E_0^{1,-1} = \ZZ/2$ on every page. The spectral sequences are isomorphic, but the abutments
$\ZZ/4$ and $\ZZ/2 \oplus \ZZ/2$ are not. Over a field this cannot happen, since every extension of vector spaces splits and $H^n \cong \bigoplus_pE_\infty^{p,n-p}$.
\end{counterexample}

\section{The spectral sequence of a filtered complex}\label{homological-algebra:section-filtered-complex}

\begin{definition}\label{homological-algebra:definition-filtered-complex}
A \emph{filtered complex} is a complex $K^\bullet$ with subcomplexes $\cdots \supseteq F^pK \supseteq F^{p+1}K \supseteq \cdots$. The
filtration is \emph{finite in each degree} if for every $n$ there are $a \leq b$ with $F^aK^n = K^n$ and $F^bK^n = 0$. The
\emph{associated graded} complexes are $\mathrm{gr}^pK = F^pK/F^{p+1}K$, and the induced filtration on cohomology is
$F^pH^n(K) = \im(H^n(F^pK) \to H^n(K))$.
\end{definition}

\begin{definition}\label{homological-algebra:definition-z-b}
For a filtered complex $K$ we abbreviate $F^p = F^pK$ and, for $r \geq 0$,
\[
Z_r^p = F^p \cap d^{-1}(F^{p+r}), \qquad B_r^p = F^p \cap d(F^{p-r+1}), \qquad
E_r^p = Z_r^p/(B_r^p + Z_{r-1}^{p+1}),
\]
with the convention $Z_{-1}^{p+1} = F^{p+1}$. These are graded by the degree of $K$, and we write $E_r^{p,q}$ for the part of $E_r^p$ in
degree $p + q$. Let $d_r \colon E_r^p \to E_r^{p+r}$ be induced by $d$.
\end{definition}

\begin{lemma}\label{homological-algebra:lemma-well-defined}
The inclusions $B_r^p \subseteq Z_r^p$ and $Z_{r-1}^{p+1} \subseteq Z_r^p$ hold, so $E_r^p$ is defined, and $d_r$ is a well-defined
morphism with $d_r^2 = 0$, of bidegree $(r, 1 - r)$. Moreover $E_0^p = \mathrm{gr}^pK$ with $d_0$ induced by $d$, and
$E_1^{p,q} = H^{p+q}(\mathrm{gr}^pK)$.
\end{lemma}

\begin{proof}
Elements of $B_r^p$ are boundaries, hence cycles; and $x \in F^{p+1}$ with $dx \in F^{p+r}$ lies in $Z_r^p$. If $x \in Z_r^p$ then
$dx \in F^{p+r}$ and $d(dx) = 0$, so $dx \in Z_r^{p+r}$. We must check that $d$ maps $B_r^p + Z_{r-1}^{p+1}$ into
$B_r^{p+r} + Z_{r-1}^{p+r+1}$: $d(B_r^p) = 0$, and for $x \in Z_{r-1}^{p+1}$ we have $dx \in F^{p+r} \cap d(F^{p+1}) = B_r^{p+r}$.
Clearly $d_r^2 = 0$. For $r = 0$: $Z_0^p = F^p$ and $B_0^p = d(F^{p+1}) \subseteq F^{p+1}$, so $E_0^p = F^p/F^{p+1}$. For $r = 1$:
$Z_1^p$ consists of the elements of $F^p$ whose image in $\mathrm{gr}^p$ is a cocycle, $Z_0^{p+1} = F^{p+1}$ and $B_1^p = d(F^p)$, so
$E_1^p = Z_1^p/(d(F^p) + F^{p+1})$ is the cohomology of $\mathrm{gr}^pK$.
\end{proof}

\begin{theorem}\label{homological-algebra:theorem-filtered-complex}
For every filtered complex, the objects $E_r^{p,q}$ and differentials $d_r$, $r \geq 0$, form a spectral sequence:
$H(E_r, d_r) \cong E_{r+1}$. If the filtration is finite in each degree, then for fixed $(p, q)$ and all sufficiently large $r$,
$E_r^{p,q} = E_\infty^{p,q}$ with
\[
E_\infty^{p,q} \cong F^pH^{p+q}(K)/F^{p+1}H^{p+q}(K).
\]
In particular, if $F^0K = K$ and $F^{n+1}K^n = 0$ for all $n$, the spectral sequence is first-quadrant (from $E_1$ on, since
$E_1^{p,q}$ vanishes for $p < 0$ or $q < 0$) and $E_1^{p,q} = H^{p+q}(\mathrm{gr}^pK) \Rightarrow H^{p+q}(K)$.
\end{theorem}

\begin{proof}
\emph{Kernel of $d_r$ at $E_r^p$.} Let $x \in Z_r^p$ with $dx \in B_r^{p+r} + Z_{r-1}^{p+r+1}$, say $dx = dy + z$ with $y \in F^{p+1}$,
$dy \in F^{p+r}$ and $z \in F^{p+r+1}$. Then $y \in Z_{r-1}^{p+1}$, and $x - y \in F^p$ with $d(x - y) = z \in F^{p+r+1}$, so
$x - y \in Z_{r+1}^p$. Conversely, $Z_{r+1}^p$ maps into the kernel, since $dx \in F^{p+r+1}$ and $d(dx) = 0$ give
$dx \in Z_{r-1}^{p+r+1}$. Hence the kernel is $(Z_{r+1}^p + Z_{r-1}^{p+1})/(B_r^p + Z_{r-1}^{p+1})$.

\emph{Image of $d_r$ in $E_r^p$.} It consists of the classes of $dx$ with $x \in Z_r^{p-r}$, that is of the elements of
$F^p \cap d(F^{p-r}) = B_{r+1}^p$. Since $B_r^p \subseteq B_{r+1}^p$, the image is $(B_{r+1}^p + Z_{r-1}^{p+1})/(B_r^p + Z_{r-1}^{p+1})$.

\emph{Cohomology.} Therefore
$H(E_r)^p = (Z_{r+1}^p + Z_{r-1}^{p+1})/(B_{r+1}^p + Z_{r-1}^{p+1})$. The map $Z_{r+1}^p \to H(E_r)^p$ is surjective, with kernel
$Z_{r+1}^p \cap (B_{r+1}^p + Z_{r-1}^{p+1}) = B_{r+1}^p + (Z_{r+1}^p \cap Z_{r-1}^{p+1})$ (as $B_{r+1}^p \subseteq Z_{r+1}^p$). Now
$Z_{r+1}^p \cap Z_{r-1}^{p+1} = F^{p+1} \cap d^{-1}(F^{p+r+1}) = Z_r^{p+1}$. So $H(E_r)^p \cong Z_{r+1}^p/(B_{r+1}^p + Z_r^{p+1}) = E_{r+1}^p$.

\emph{Convergence.} Fix $n$ and $p$ and let $Z$, $B$ denote the cocycles and coboundaries of $K$. If the filtration is finite in degrees
$n - 1$, $n$ and $n + 1$, then for large $r$ we have $F^{p+r}K^{n+1} = 0$ and $F^{p-r+1}K^{n-1} = K^{n-1}$, so in degree $n$:
$Z_r^p = F^p \cap Z$, $Z_{r-1}^{p+1} = F^{p+1} \cap Z$ and $B_r^p = F^p \cap B$. Hence
$E_\infty^p = (F^p \cap Z)/((F^p \cap B) + (F^{p+1} \cap Z))$. On the other hand $F^pH^n = (F^p \cap Z + B)/B$, and
\[
\frac{F^pH^n}{F^{p+1}H^n} \cong \frac{F^p \cap Z}{(F^p \cap Z) \cap (F^{p+1} \cap Z + B)} = \frac{F^p \cap Z}{F^{p+1} \cap Z + F^p \cap B},
\]
using the modular law. For the last statement, $\mathrm{gr}^pK = 0$ for $p < 0$ and $\mathrm{gr}^pK^n = 0$ for $p > n$, so $E_1^{p,q} = 0$ unless
$p \geq 0$ and $q \geq 0$, and $F^pH^n$ runs from $H^n$ to $0$.
\end{proof}

\section{Double complexes}\label{homological-algebra:section-double-complexes-spectral-sequences}

\begin{theorem}\label{homological-algebra:theorem-double-complex}
Let $C$ be a first-quadrant double complex (Definition~\ref{homological-algebra:definition-double-complex}) and $T = \mathrm{Tot}(C)$.
There are two convergent first-quadrant spectral sequences
\[
{}^{I}E_2^{p,q} = H^p_h(H^q_v(C)) \Rightarrow H^{p+q}(T), \qquad {}^{II}E_2^{p,q} = H^p_v(H^q_h(C)) \Rightarrow H^{p+q}(T),
\]
where $H_v^q(C)$ is the cohomology of the columns (a complex in $p$ via $d_h$), and $H_h^q(C)$ that of the rows (a complex in the
other index via $d_v$).
\end{theorem}

\begin{proof}
Filter $T$ by columns: ${}^IF^pT^n = \bigoplus_{i \geq p}C^{i, n-i}$. These are subcomplexes, $F^0T = T$ and $F^{n+1}T^n = 0$. The graded
piece $\mathrm{gr}^pT$ is the column $C^{p,\bullet}$ (shifted by $p$) with differential $\pm d_v$, so
${}^IE_1^{p,q} = H^q_v(C^{p,\bullet})$ by Theorem~\ref{homological-algebra:theorem-filtered-complex}. The differential $d_1$ is the
connecting map of $0 \to \mathrm{gr}^{p+1} \to F^p/F^{p+2} \to \mathrm{gr}^p \to 0$; a cocycle $x \in C^{p,q}$ of $d_v$ lifts to itself, and its total
differential is $d_hx \in C^{p+1,q}$, so $d_1 = H^q_v(d_h)$ and ${}^IE_2^{p,q} = H^p_h(H^q_v(C))$. Convergence follows from
Theorem~\ref{homological-algebra:theorem-filtered-complex}. Filtering by rows, ${}^{II}F^pT^n = \bigoplus_{j \geq p}C^{n-j, j}$, gives the other
spectral sequence in the same way.
\end{proof}

\begin{example}\label{homological-algebra:example-acyclic-assembly}
If the rows of $C$ are exact in positive degrees, with kernel complex $K^\bullet$ in column $0$, then ${}^{II}E_1^{p,q} = H^q_h(C^{\bullet,p})$
vanishes for $q > 0$ and equals $K^p$ for $q = 0$. So ${}^{II}E_2^{p,0} = H^p(K)$, all other terms vanish, and
$H^n(T) \cong H^n(K)$ by Proposition~\ref{homological-algebra:proposition-edge-maps}(3). This recovers
Lemma~\ref{homological-algebra:lemma-acyclic-assembly}.
\end{example}

\section{Exact couples}\label{homological-algebra:section-exact-couples}

\begin{definition}\label{homological-algebra:definition-exact-couple}
An \emph{exact couple} consists of objects $D$, $E$ and morphisms $i \colon D \to D$, $j \colon D \to E$, $k \colon E \to D$ such that the triangle
$D \xrightarrow{i} D \xrightarrow{j} E \xrightarrow{k} D$ is exact at each vertex. Put $d = jk \colon E \to E$; then $d^2 = j(kj)k = 0$.
\end{definition}

\begin{proposition}\label{homological-algebra:proposition-derived-couple}
Given an exact couple, let $D' = i(D)$, $E' = \ker d/\im d$, $i' = i|_{D'}$, $j'(i(x)) = [j(x)]$ and $k'([e]) = k(e)$. Then $(D', E', i', j', k')$
is an exact couple, the \emph{derived couple}. Iterating gives a spectral sequence $E^{(r)}$ with differentials $d^{(r)} = j^{(r)}k^{(r)}$.
\end{proposition}

\begin{proof}
$j'$ is well defined: if $i(x) = i(x')$ then $x - x' = k(e)$ and $j(x - x') = jk(e) = d(e)$; and $j(x)$ is a $d$-cocycle since
$dj(x) = jkj(x) = 0$. $k'$ is well defined: if $e \in \ker d$ then $jk(e) = 0$ so $k(e) \in \ker j = \im i = D'$, and $k(d(e')) = kjk(e') = 0$.
Exactness at $D'$ (between $i'$ and $j'$): $j'(i(x)) = 0$ iff $j(x) = jk(e)$ iff $x - k(e) \in \ker j = \im i$, iff $i(x) \in i(\im i) + ik(E) = i(D')$
since $ik = 0$. Exactness at $E'$: $k'[e] = 0$ iff $k(e) = 0$ iff $e = j(x)$, iff $[e] = j'(i(x))$. Exactness at the other $D'$
(between $k'$ and $i'$): $i'(i(x)) = 0$ iff $i(x) \in \ker i = \im k$, iff $i(x) = k(e)$ for some $e$, which is then a $d$-cocycle since
$jk(e) = ji(x) = 0$; that is, iff $i(x) \in \im k'$.
\end{proof}

\begin{remark}\label{homological-algebra:remark-exact-couple-filtered}
A filtered complex gives an exact couple with $D = \bigoplus_{p,n}H^n(F^pK)$, $E = \bigoplus_{p,n}H^n(\mathrm{gr}^pK)$, $i$ induced by the inclusions
$F^{p+1} \subseteq F^p$, and $j$, $k$ from the long exact sequences of $0 \to F^{p+1} \to F^p \to \mathrm{gr}^p \to 0$. The resulting spectral sequence is
the one of Theorem~\ref{homological-algebra:theorem-filtered-complex}, from $E_1$ on \cite[Section 5.9]{weibel-homological}.
\end{remark}

\section{The Grothendieck spectral sequence}\label{homological-algebra:section-grothendieck}

\begin{definition}\label{homological-algebra:definition-cartan-eilenberg}
Let $A^\bullet$ be a complex with $A^n = 0$ for $n < 0$ in an abelian category with enough injectives. A \emph{Cartan--Eilenberg resolution}
of $A$ is a first-quadrant double complex $I^{\bullet,\bullet}$ of injectives with a morphism $A^p \to I^{p,0}$ for each $p$, such that for every $p$ the
columns give injective resolutions $A^p \to I^{p,\bullet}$, $Z^p(A) \to Z^p_h(I)^\bullet$, $B^p(A) \to B^p_h(I)^\bullet$ and $H^p(A) \to H^p_h(I)^\bullet$, where
$Z_h$, $B_h$, $H_h$ denote horizontal cocycles, coboundaries and cohomology.
\end{definition}

\begin{lemma}\label{homological-algebra:lemma-cartan-eilenberg-exists}
Every such complex $A$ has a Cartan--Eilenberg resolution.
\end{lemma}

\begin{proof}
For each $p$ choose injective resolutions $B^p(A) \to J_B^{p,\bullet}$ and $H^p(A) \to J_H^{p,\bullet}$. The horseshoe lemma
(Lemma~\ref{homological-algebra:lemma-horseshoe}) applied to $0 \to B^p \to Z^p \to H^p \to 0$ gives a resolution $J_Z^{p,\bullet}$ of $Z^p$ with
$J_Z^{p,q} = J_B^{p,q} \oplus J_H^{p,q}$, and applied to $0 \to Z^p \to A^p \to B^{p+1} \to 0$ gives a resolution $I^{p,\bullet}$ of $A^p$ with
$I^{p,q} = J_Z^{p,q} \oplus J_B^{p+1,q}$. Define $d_h \colon I^{p,q} \to I^{p+1,q}$ as the composite
$I^{p,q} \to J_B^{p+1,q} \to J_Z^{p+1,q} \to I^{p+1,q}$ of the projection and inclusions provided by the horseshoe lemma. Then $d_h^2 = 0$ (the
projection $I^{p+1, q} \to J_B^{p+2,q}$ kills $J_Z^{p+1,q}$), $d_h$ commutes with the vertical differentials since all maps involved are morphisms
of complexes, and it covers $d \colon A^p \to A^{p+1}$. By construction the horizontal cocycles, coboundaries and cohomology in column $p$ are
$J_Z^{p,\bullet}$, $J_B^{p,\bullet}$ and $J_H^{p,\bullet}$.
\end{proof}

\begin{theorem}[Grothendieck spectral sequence]\label{homological-algebra:theorem-grothendieck}
Let $\mathcal{A}, \mathcal{B}, \mathcal{C}$ be abelian categories, $\mathcal{A}$ and $\mathcal{B}$ with enough injectives, and $F \colon \mathcal{A} \to \mathcal{B}$,
$G \colon \mathcal{B} \to \mathcal{C}$ left exact functors such that $F$ maps injective objects to $G$-acyclic objects. For every $X \in \mathcal{A}$ there is a
convergent first-quadrant spectral sequence
\[
E_2^{p,q} = R^pG(R^qF(X)) \Rightarrow R^{p+q}(G \circ F)(X).
\]
\end{theorem}

\begin{proof}
Let $X \to A^\bullet$ be an injective resolution in $\mathcal{A}$, and let $I^{\bullet,\bullet}$ be a Cartan--Eilenberg resolution of the complex $F(A^\bullet)$ in
$\mathcal{B}$. Consider the double complex $G(I)$ and its total complex $T$.

\emph{First filtration (by columns).} For each $p$, the column $I^{p,\bullet}$ is an injective resolution of the $G$-acyclic object $F(A^p)$, so
$H^q_v(G(I^{p,\bullet})) = R^qG(F(A^p))$ is $G(F(A^p))$ for $q = 0$ and $0$ for $q > 0$. So ${}^IE_2^{p,q} = 0$ for $q > 0$ and
${}^IE_2^{p,0} = H^p(GF(A^\bullet)) = R^p(GF)(X)$, and $H^n(T) \cong R^n(GF)(X)$ by Proposition~\ref{homological-algebra:proposition-edge-maps}(3).

\emph{Second filtration (by rows).} Fix a row index $j$. In the row $I^{\bullet,j}$ the sequences
$0 \to B^{q,j}_h \to Z^{q,j}_h \to H^{q,j}_h \to 0$ and $0 \to Z^{q,j}_h \to I^{q,j} \to B^{q+1,j}_h \to 0$ consist of injective objects (they are terms of
injective resolutions), hence split; so $G$ commutes with taking horizontal cocycles, coboundaries and cohomology of this row, and
$H^q_h(G(I^{\bullet,j})) = G(H^{q,j}_h(I))$. For fixed $q$, the column $H^{q,\bullet}_h(I)$ is an injective resolution of
$H^q(F(A^\bullet)) = R^qF(X)$. By Theorem~\ref{homological-algebra:theorem-double-complex},
${}^{II}E_2^{p,q} = H^p_v(H^q_h(G(I))) = H^p(G(H^{q,\bullet}_h(I))) = R^pG(R^qF(X))$, and it converges to
$H^{p+q}(T) \cong R^{p+q}(GF)(X)$.
\end{proof}

\section{Derived Categories}\label{homological-algebra:section-derived-categories}
The derived category of an abelian category is obtained from complexes by declaring quasi-isomorphisms to be
isomorphisms. It is the natural home of derived functors: instead of a sequence of functors $R^nF$ one has a
single functor $RF$ on complexes, compositions of derived functors are derived functors of compositions, and
hypercohomology of complexes (for instance of the de Rham complex of a manifold) is defined without choices.
Filtered derived categories and strict morphisms, discussed at the end, are the framework of Deligne's mixed Hodge
theory (Chapter~\ref{mixed-hodge-structures}).

Some foundational verifications (the octahedral axiom, and the construction of localisations by roofs) are quoted
from \cite[Chapter 10]{weibel-homological}, \cite[Chapter III]{gelfand-manin} and \cite{kashiwara-schapira}; everything
used about derived functors is proved.

\noindent\textbf{Prerequisites.} Chapter~\ref{categories} (Categories), Chapter~\ref{homological-algebra} (Homological
Algebra) and Section~\ref{homological-algebra:section-spectral-sequences} (Spectral Sequences).

Throughout, $\mathcal{A}$ is an abelian category, complexes are as in
Definition~\ref{homological-algebra:definition-complex}, and we argue with elements as justified by
Remark~\ref{homological-algebra:remark-freyd-mitchell}.

\section{The homotopy category}\label{homological-algebra:section-homotopy-category}

\begin{definition}\label{homological-algebra:definition-homotopy-category}
The \emph{homotopy category} $K(\mathcal{A})$ has the complexes as objects and homotopy classes of morphisms of complexes
(Definition~\ref{homological-algebra:definition-homotopy}) as morphisms. $K^+(\mathcal{A})$ is its full subcategory of bounded
below complexes. The \emph{shift} $X \mapsto X[1]$ is an automorphism of $K(\mathcal{A})$. A \emph{distinguished triangle} in
$K(\mathcal{A})$ is a diagram $X \to Y \to Z \to X[1]$ isomorphic (in $K(\mathcal{A})$) to one of the form
$X \xrightarrow{f} Y \xrightarrow{i} \mathrm{Cone}(f) \xrightarrow{p} X[1]$, with $i$ and $p$ the maps of
Proposition~\ref{homological-algebra:proposition-cone}.
\end{definition}

\begin{remark}\label{homological-algebra:remark-homotopy-well-defined}
Homotopy is compatible with sums and compositions, so $K(\mathcal{A})$ is an additive category, and cohomology
$H^n \colon K(\mathcal{A}) \to \mathcal{A}$ is a well-defined functor by Lemma~\ref{homological-algebra:lemma-homotopy-invariance}.
\end{remark}

\begin{definition}\label{homological-algebra:definition-triangulated}
A \emph{triangulated category} is an additive category $\mathcal{T}$ with an additive automorphism $X \mapsto X[1]$ and a class of
\emph{distinguished triangles} $X \to Y \to Z \to X[1]$ such that:
\begin{enumerate}
\item[(TR1)] $X \xrightarrow{\id} X \to 0 \to X[1]$ is distinguished; triangles isomorphic to distinguished ones are distinguished;
every morphism $X \to Y$ is the first morphism of a distinguished triangle.
\item[(TR2)] $X \xrightarrow{f} Y \xrightarrow{g} Z \xrightarrow{h} X[1]$ is distinguished if and only if
$Y \xrightarrow{g} Z \xrightarrow{h} X[1] \xrightarrow{-f[1]} Y[1]$ is.
\item[(TR3)] Given distinguished triangles $X \to Y \to Z \to X[1]$ and $X' \to Y' \to Z' \to X'[1]$ and morphisms $u \colon X \to X'$,
$v \colon Y \to Y'$ making the first square commute, there is $w \colon Z \to Z'$ making the whole diagram commute.
\item[(TR4)] (Octahedral axiom) For composable $X \xrightarrow{f} Y \xrightarrow{g} Z$ and distinguished triangles on $f$, $g$ and $gf$ with
third vertices $Z'$, $X'$, $Y'$, there is a distinguished triangle $Z' \to Y' \to X' \to Z'[1]$ compatible with the given
morphisms.
\end{enumerate}
An additive functor between triangulated categories commuting with the shifts and preserving distinguished triangles is
\emph{exact} (or triangulated).
\end{definition}

\begin{theorem}\label{homological-algebra:theorem-homotopy-triangulated}
$K(\mathcal{A})$ and $K^+(\mathcal{A})$ are triangulated categories.
\end{theorem}

\begin{proof}[Sketch of proof]
(TR1): $\mathrm{Cone}(\id_X)$ is contractible, with contracting homotopy $(a, b) \mapsto (b, 0)$, so it is isomorphic to $0$ in
$K(\mathcal{A})$; the other two parts hold by definition. (TR3): if $h$ is a homotopy with $vf - f'u = dh + hd$, then
$(a, b) \mapsto (u(a), v(b) + h(a))$ is a morphism of complexes $\mathrm{Cone}(f) \to \mathrm{Cone}(f')$ (a direct check using the formula for
the differential of the cone in Definition~\ref{homological-algebra:definition-cone}) compatible with the maps of the triangles. The
omitted verifications are (TR2), which amounts to an explicit homotopy equivalence $\mathrm{Cone}(i) \simeq X[1]$ for
$i \colon Y \to \mathrm{Cone}(f)$, and the octahedral axiom (TR4); see \cite[Sections 10.1--10.2]{weibel-homological}.
\end{proof}

\begin{proposition}\label{homological-algebra:proposition-cohomological}
Let $X \xrightarrow{f} Y \xrightarrow{g} Z \xrightarrow{h} X[1]$ be a distinguished triangle in a triangulated category $\mathcal{T}$ and $W$ an
object. Then $gf = 0$, and the sequence $\Hom(W, X) \to \Hom(W, Y) \to \Hom(W, Z)$ is exact; so by (TR2) there is a long exact
sequence $\cdots \to \Hom(W, X) \to \Hom(W, Y) \to \Hom(W, Z) \to \Hom(W, X[1]) \to \cdots$. Similarly for $\Hom(-, W)$. For
$\mathcal{T} = K(\mathcal{A})$, the functor $H^0$ sends distinguished triangles to long exact sequences.
\end{proposition}

\begin{proof}
Apply (TR3) to the triangles $X \xrightarrow{\id} X \to 0 \to X[1]$ and $X \to Y \to Z \to X[1]$ with $u = \id$, $v = f$: the resulting $w$
satisfies $gf = w \circ 0 = 0$. Let $\varphi \colon W \to Y$ with $g\varphi = 0$. By (TR1) and (TR2), the rotated triangles
$W \to 0 \to W[1] \xrightarrow{-\id} W[1]$ and $Y \xrightarrow{g} Z \xrightarrow{h} X[1] \xrightarrow{-f[1]} Y[1]$ are distinguished. The square formed by
$\varphi \colon W \to Y$ and $0 \colon 0 \to Z$ commutes because $g\varphi = 0$, so (TR3) gives $w \colon W[1] \to X[1]$ with
$(-f[1]) \circ w = \varphi[1] \circ (-\id)$. Then $\psi = w[-1]$ satisfies $f\psi = \varphi$. The dual statement is proved in the same way. For $K(\mathcal{A})$ the last claim
follows from the long exact sequence of the cone (Proposition~\ref{homological-algebra:proposition-cone}) and invariance under
isomorphism.
\end{proof}

\section{The derived category}\label{homological-algebra:section-derived-category}

\begin{definition}\label{homological-algebra:definition-multiplicative-system}
A class $S$ of morphisms in a category $\mathcal{C}$ is a \emph{multiplicative system} if it contains the identities and is closed
under composition, and:
\begin{enumerate}
\item for $s \colon X' \to X$ in $S$ and $f \colon Y \to X$ there are $t \colon Y' \to Y$ in $S$ and $g \colon Y' \to X'$ with $ft = sg$, and dually;
\item if $f, g \colon X \to Y$ and $s \colon Y \to Y'$ in $S$ satisfy $sf = sg$, then $ft = gt$ for some $t \colon X' \to X$ in $S$, and dually.
\end{enumerate}
\end{definition}

\begin{theorem}[Localisation]\label{homological-algebra:theorem-localisation}
Let $S$ be a multiplicative system in $\mathcal{C}$. There is a category $\mathcal{C}[S^{-1}]$ with the same objects and a functor
$Q \colon \mathcal{C} \to \mathcal{C}[S^{-1}]$ that is the identity on objects, sends $S$ to isomorphisms, and is universal: every functor sending
$S$ to isomorphisms factors uniquely through $Q$. Morphisms $X \to Y$ in $\mathcal{C}[S^{-1}]$ are equivalence classes of \emph{roofs}
$X \xleftarrow{s} X' \xrightarrow{f} Y$ with $s \in S$, representing $Q(f)Q(s)^{-1}$, and for $f, g \colon X \to Y$ in $\mathcal{C}$ one has
$Q(f) = Q(g)$ if and only if $ft = gt$ for some $t \in S$. If $\mathcal{C}$ is additive (triangulated) and $S$ is compatible
with the structure, then $\mathcal{C}[S^{-1}]$ is additive (triangulated), with distinguished triangles the images of those of $\mathcal{C}$.
\end{theorem}

This is proved in \cite[Chapter III]{gelfand-manin} and \cite[Chapter 7]{kashiwara-schapira}; for the triangulated case see
\cite[Section 10.3]{weibel-homological}.

\begin{proposition}\label{homological-algebra:proposition-qis-multiplicative}
The quasi-isomorphisms form a multiplicative system in $K(\mathcal{A})$ and in $K^+(\mathcal{A})$.
\end{proposition}

\begin{proof}
Identities and composites of quasi-isomorphisms are quasi-isomorphisms. Note that a morphism is a quasi-isomorphism iff its cone
is acyclic (Proposition~\ref{homological-algebra:proposition-cone}), and by
Proposition~\ref{homological-algebra:proposition-cohomological} the cohomology of the three vertices of a distinguished triangle
forms a long exact sequence.

(1) Let $s \colon X' \to X$ be a quasi-isomorphism and $f \colon Y \to X$. Complete $s$ to a triangle $X' \to X \xrightarrow{p} C \to X'[1]$, so $C$ is
acyclic, and complete $pf \colon Y \to C$ to a triangle $Y' \xrightarrow{t} Y \xrightarrow{pf} C \to Y'[1]$ (using (TR1) and (TR2)). Then $t$ is a
quasi-isomorphism since $C$ is acyclic, and (TR3) applied to the rotated triangles gives $g \colon Y' \to X'$ with $sg = ft$. The dual
statement is proved in the same way.

(2) Suppose $s(f - g) = 0$ with $s \colon Y \to Y'$ a quasi-isomorphism. Complete $s$ to a triangle $C[-1] \xrightarrow{k} Y \xrightarrow{s} Y' \to C$ with $C$
acyclic. By Proposition~\ref{homological-algebra:proposition-cohomological}, $f - g = k\varphi$ for some $\varphi \colon X \to C[-1]$. Complete $\varphi$ to a
triangle $X' \xrightarrow{t} X \xrightarrow{\varphi} C[-1] \to X'[1]$; then $t$ is a quasi-isomorphism and $(f - g)t = k\varphi t = 0$. Dually.
\end{proof}

\begin{definition}\label{homological-algebra:definition-derived-category}
The \emph{derived category} is $D(\mathcal{A}) = K(\mathcal{A})[\mathrm{Qis}^{-1}]$, and $D^+(\mathcal{A}) = K^+(\mathcal{A})[\mathrm{Qis}^{-1}]$. By
Theorem~\ref{homological-algebra:theorem-localisation} and Proposition~\ref{homological-algebra:proposition-qis-multiplicative} they are
triangulated, $H^n$ factors through them, and a morphism in $D(\mathcal{A})$ is an isomorphism if and only if it induces
isomorphisms on all $H^n$.
\end{definition}

\begin{remark}\label{homological-algebra:remark-size}
The morphism classes of $D(\mathcal{A})$ need not be small sets for an arbitrary locally small $\mathcal{A}$, since roofs involve all complexes.
For $\mathcal{A}$ with enough injectives this problem disappears for $D^+(\mathcal{A})$ by
Corollary~\ref{homological-algebra:corollary-hom-injectives}, which computes morphisms using a single complex of injectives.
\end{remark}

\section{Injective resolutions of complexes}\label{homological-algebra:section-injective-resolutions}

\begin{lemma}\label{homological-algebra:lemma-acyclic-to-injective}
Let $X$ be an acyclic complex and $I$ a bounded below complex of injective objects. Every morphism of complexes $f \colon X \to I$ is
homotopic to zero.
\end{lemma}

\begin{proof}
Let $I^n = 0$ for $n < n_0$; put $h^n = 0 \colon X^n \to I^{n-1}$ for $n \leq n_0$. Suppose $h^k$ is defined for $k \leq n$ with
$f^{k-1} = dh^{k-1} + h^kd$ for $k \leq n$. Let $g = f^n - dh^n \colon X^n \to I^n$. Then
$gd = df^{n-1} - dh^nd = df^{n-1} - d(f^{n-1} - dh^{n-1}) = 0$, so $g$ vanishes on $\im(X^{n-1} \to X^n) = \ker(X^n \to X^{n+1})$. Hence $g$
factors through $X^n/\ker \cong \im(X^n \to X^{n+1}) \subseteq X^{n+1}$ and extends, by injectivity of $I^n$, to $h^{n+1} \colon X^{n+1} \to I^n$ with
$h^{n+1}d = g$, that is, $f^n = dh^n + h^{n+1}d$.
\end{proof}

\begin{corollary}\label{homological-algebra:corollary-hom-injectives}
Let $I$ be a bounded below complex of injectives. For every quasi-isomorphism $s \colon X' \to X$, composition with $s$ is a bijection
$\Hom_{K(\mathcal{A})}(X, I) \to \Hom_{K(\mathcal{A})}(X', I)$. Consequently $\Hom_{D(\mathcal{A})}(X, I) = \Hom_{K(\mathcal{A})}(X, I)$, and every quasi-isomorphism
$I \to I'$ between bounded below complexes of injectives is a homotopy equivalence.
\end{corollary}

\begin{proof}
Complete $s$ to a distinguished triangle $X' \to X \to C \to X'[1]$ with $C$ acyclic. By Lemma~\ref{homological-algebra:lemma-acyclic-to-injective},
$\Hom_K(C, I) = \Hom_K(C[-1], I) = 0$ ($C[-1]$ is also acyclic), and the long exact sequence of
Proposition~\ref{homological-algebra:proposition-cohomological} gives the bijection. For the second statement, we use the description of morphisms in
Theorem~\ref{homological-algebra:theorem-localisation}. The map $Q \colon \Hom_K(X, I) \to \Hom_D(X, I)$ is surjective: a roof
$X \xleftarrow{s} X' \xrightarrow{f} I$ has $f = f's$ for some $f' \colon X \to I$, and then $Q(f)Q(s)^{-1} = Q(f')$. It is injective: if $Q(f_1) = Q(f_2)$,
then $f_1t = f_2t$ for a quasi-isomorphism $t$, so $f_1 = f_2$ by the bijection just proved. For the last, a quasi-isomorphism $I \to I'$ is an isomorphism in
$D(\mathcal{A})$, hence by the second statement an isomorphism in $K(\mathcal{A})$.
\end{proof}

\begin{proposition}\label{homological-algebra:proposition-resolutions-complexes}
If $\mathcal{A}$ has enough injectives, every bounded below complex $X$ admits a quasi-isomorphism $X \to I$ to a bounded below complex of
injectives. Hence $K^+(\mathrm{Inj}\,\mathcal{A}) \to D^+(\mathcal{A})$ is an equivalence of categories.
\end{proposition}

\begin{proof}
Shifting, assume $X^n = 0$ for $n < 0$. Let $J^{\bullet,\bullet}$ be a Cartan--Eilenberg resolution of $X$
(Lemma~\ref{homological-algebra:lemma-cartan-eilenberg-exists}) and $I = \mathrm{Tot}(J)$, a bounded below complex of injectives (each $I^n$ is a
finite sum). The augmentations $X^p \to J^{p,0}$ define a morphism $X \to I$. For each $p$ the column $0 \to X^p \to J^{p,0} \to J^{p,1} \to \cdots$ is
exact, so by Lemma~\ref{homological-algebra:lemma-acyclic-assembly} (applied with the roles of rows and columns exchanged) the map
$X \to I$ is a quasi-isomorphism. (The total complex of the transposed double complex is isomorphic to $I$ by the signs $(-1)^{pq}$, and the
compatibility of $X^p \to J^{p,0}$ with the differentials holds because the column differential kills the image of $X^p$.) The functor $K^+(\mathrm{Inj}\,\mathcal{A}) \to D^+(\mathcal{A})$ is fully faithful by
Corollary~\ref{homological-algebra:corollary-hom-injectives} and essentially surjective by the first statement.
\end{proof}

\begin{definition}\label{homological-algebra:definition-right-derived}
Let $\mathcal{A}$ have enough injectives and $F \colon \mathcal{A} \to \mathcal{B}$ be an additive functor to an abelian category. For each bounded
below complex $X$ choose a quasi-isomorphism $X \to I_X$ as in Proposition~\ref{homological-algebra:proposition-resolutions-complexes}. The
\emph{right derived functor} is $RF \colon D^+(\mathcal{A}) \to D^+(\mathcal{B})$, $RF(X) = F(I_X)$ (with $F$ applied termwise), on morphisms via
$\Hom_{D}(X, Y) \cong \Hom_K(I_X, I_Y)$. For $X \in D^+(\mathcal{A})$ and $n \in \ZZ$, $R^nF(X) = H^n(RF(X))$ is the $n$-th \emph{hypercohomology}
of $F$ at $X$.
\end{definition}

\begin{proposition}\label{homological-algebra:proposition-derived-functor-properties}
$RF$ is an exact functor of triangulated categories, independent of the choices up to canonical isomorphism. If $F$ is left exact and
$X \in \mathcal{A}$ is viewed as a complex in degree $0$, then $H^n(RF(X)) = R^nF(X)$ as in
Definition~\ref{homological-algebra:definition-derived-functor}. For $X, Y \in \mathcal{A}$ and $n \in \ZZ$,
$\Hom_{D(\mathcal{A})}(X, Y[n]) \cong R^n\Hom(X, -)(Y)$; in particular $\Hom_{D(R\text{-}\cat{Mod})}(M, N[n]) \cong \Ext^n_R(M, N)$ for
$R$-modules $M$, $N$.
\end{proposition}

\begin{proof}
Independence of the choices follows from uniqueness of $I_X$ up to unique isomorphism in $K(\mathcal{A})$
(Corollary~\ref{homological-algebra:corollary-hom-injectives}). For exactness, a distinguished triangle of $D^+(\mathcal{A})$ is isomorphic to one
of the form $X \xrightarrow{g} Y \to \mathrm{Cone}(g) \to X[1]$ with $g$ a morphism of bounded below complexes. By
Corollary~\ref{homological-algebra:corollary-hom-injectives} there is $g_I \colon I_X \to I_Y$ with $g_I\iota_X = \iota_Y g$ in $K(\mathcal{A})$, where
$\iota$ denotes the resolutions, and (TR3) in $K(\mathcal{A})$ gives a morphism of triangles $(\iota_X, \iota_Y, w)$ from the triangle of $g$ to that of
$g_I$. By Proposition~\ref{homological-algebra:proposition-cohomological}, the five lemma (Lemma~\ref{homological-algebra:lemma-five}) shows
that $w$ is a quasi-isomorphism. So the triangle is isomorphic in $D(\mathcal{A})$ to $I_X \to I_Y \to \mathrm{Cone}(g_I) \to I_X[1]$, whose vertices are
bounded below complexes of injectives, and $RF$ sends it to $F(I_X) \to F(I_Y) \to F(\mathrm{Cone}(g_I)) \to F(I_X)[1]$. Since $F$ is additive and applied
termwise, $F(\mathrm{Cone}(g_I)) = \mathrm{Cone}(F(g_I))$ with the corresponding maps, so this triangle is distinguished.

For $X \in \mathcal{A}$, an injective resolution is a quasi-isomorphism $X \to I$, so $H^n(RF(X)) = H^n(F(I)) = R^nF(X)$. Finally
$\Hom_D(X, Y[n]) = \Hom_K(X, I_Y[n])$ by Corollary~\ref{homological-algebra:corollary-hom-injectives}. A morphism of complexes $X \to I_Y[n]$ is a
morphism $\varphi \colon X \to I_Y^n$ with $d\varphi = 0$, that is, an element of $\ker(\Hom(X, I_Y^n) \to \Hom(X, I_Y^{n+1}))$, and it is homotopic to
zero if and only if $\varphi = dh$ for some $h \colon X \to I_Y^{n-1}$. So $\Hom_K(X, I_Y[n]) = H^n(\Hom(X, I_Y)) = R^n\Hom(X, -)(Y)$. For modules this
is $\Ext^n_R(M, N)$ by Definition~\ref{homological-algebra:definition-ext-tor}, as $R\text{-}\cat{Mod}$ has enough injectives
(Theorem~\ref{homological-algebra:theorem-enough-injectives}).
\end{proof}

\begin{lemma}\label{homological-algebra:lemma-acyclic-complexes}
Let $\mathcal{B}$ have enough injectives, $G \colon \mathcal{B} \to \mathcal{C}$ be left exact, and $Y$ a bounded below acyclic complex of
$G$-acyclic objects (Proposition~\ref{homological-algebra:proposition-acyclic-resolutions}). Then $G(Y)$ is acyclic.
Consequently, if $Y$ is any bounded below complex of $G$-acyclic objects, then $RG(Y) \cong G(Y)$ in $D^+(\mathcal{C})$.
\end{lemma}

\begin{proof}
Let $Z^n = \ker(Y^n \to Y^{n+1})$. Acyclicity gives short exact sequences $0 \to Z^n \to Y^n \to Z^{n+1} \to 0$, and $Z^n = 0$ for $n \ll 0$. By induction
and the long exact sequence of Theorem~\ref{homological-algebra:theorem-derived-functors}, all $Z^n$ are $G$-acyclic, so $G$ preserves the
exactness of these sequences ($R^1G(Z^n) = 0$). Since $G(Z^{n+1}) \to G(Y^{n+1})$ is a monomorphism, the kernel of $G(Y^n) \to G(Y^{n+1})$ is
$G(Z^n)$, which is the image of $G(Y^{n-1}) \to G(Z^n) \to G(Y^n)$; so $G(Y)$ is acyclic. For the second statement, let $Y \to I$ be a quasi-isomorphism to a complex of
injectives; its cone is a bounded below acyclic complex (Proposition~\ref{homological-algebra:proposition-cone}) of $G$-acyclic objects
(termwise $Y^{n+1} \oplus I^n$, and $R^kG$ is additive). By the first statement $G(\mathrm{Cone}) = \mathrm{Cone}(G(Y) \to G(I))$ is acyclic, so
$G(Y) \to G(I) = RG(Y)$ is a quasi-isomorphism.
\end{proof}

\begin{theorem}\label{homological-algebra:theorem-composition}
Let $F \colon \mathcal{A} \to \mathcal{B}$ and $G \colon \mathcal{B} \to \mathcal{C}$ be left exact, with $\mathcal{A}$ and $\mathcal{B}$ having enough injectives,
and suppose $F$ sends injectives to $G$-acyclic objects. Then $R(G \circ F) \cong RG \circ RF$.
\end{theorem}

\begin{proof}
For $X \in D^+(\mathcal{A})$ with resolution $X \to I_X$, $RF(X) = F(I_X)$ is a bounded below complex of $G$-acyclic objects, so
$RG(RF(X)) \cong G(F(I_X)) = R(GF)(X)$ by Lemma~\ref{homological-algebra:lemma-acyclic-complexes}.
\end{proof}

\begin{proposition}[Hypercohomology spectral sequences]\label{homological-algebra:proposition-hypercohomology}
Let $F$ be left exact, $\mathcal{A}$ with enough injectives, and $X$ a complex with $X^n = 0$ for $n < 0$. There are convergent first-quadrant
spectral sequences
\[
{}'E_1^{p,q} = R^qF(X^p) \Rightarrow R^{p+q}F(X), \qquad {}''E_2^{p,q} = R^pF(H^q(X)) \Rightarrow R^{p+q}F(X).
\]
\end{proposition}

\begin{proof}
Let $J$ be a Cartan--Eilenberg resolution of $X$; by the proof of Proposition~\ref{homological-algebra:proposition-resolutions-complexes},
$RF(X) = F(\mathrm{Tot}(J)) = \mathrm{Tot}(F(J))$. Filtering by columns as in the proof of Theorem~\ref{homological-algebra:theorem-double-complex},
${}'E_1^{p,q} = H^q_v(F(J^{p,\bullet})) = R^qF(X^p)$, since $J^{p,\bullet}$ is an injective resolution of $X^p$. Filtering by rows, the rows of $J$ are split
(second filtration in the proof of Theorem~\ref{homological-algebra:theorem-grothendieck}), so $H^q_h(F(J)) = F(H^q_h(J))$, and $H^q_h(J)^{\bullet}$ is
an injective resolution of $H^q(X)$; hence ${}''E_2^{p,q} = H^p(F(H^q_h(J))) = R^pF(H^q(X))$. Convergence is
Theorem~\ref{homological-algebra:theorem-double-complex}.
\end{proof}

\section{$t$-structures}\label{homological-algebra:section-t-structures}

\begin{definition}\label{homological-algebra:definition-truncation}
For a complex $X$ and $n \in \ZZ$ the \emph{truncations} are
\[
\tau_{\leq n}X = (\cdots \to X^{n-1} \to Z^n(X) \to 0 \to \cdots), \qquad
\tau_{\geq n}X = (\cdots \to 0 \to X^n/B^n(X) \to X^{n+1} \to \cdots).
\]
\end{definition}

\begin{proposition}\label{homological-algebra:proposition-truncation}
There are natural morphisms $\tau_{\leq n}X \to X \to \tau_{\geq n}X$, and $H^k(\tau_{\leq n}X) = H^k(X)$ for $k \leq n$ and $0$ otherwise, while
$H^k(\tau_{\geq n}X) = H^k(X)$ for $k \geq n$ and $0$ otherwise. There is a distinguished triangle
$\tau_{\leq n}X \to X \to \tau_{\geq n+1}X \to (\tau_{\leq n}X)[1]$ in $D(\mathcal{A})$.
\end{proposition}

\begin{proof}
The cohomology computations are immediate from the definitions. The morphism $X \to \tau_{\geq n+1}X$ is surjective with kernel the complex
$K = (\cdots \to X^{n-1} \to X^n \to B^{n+1}(X) \to 0 \to \cdots)$, and $\tau_{\leq n}X \to K$ is a quasi-isomorphism (both have cohomology $H^k(X)$ for
$k \leq n$ and zero above, $H^{n+1}(K) = 0$ as $X^n \to B^{n+1}$ is onto). Finally, a short exact sequence of complexes $0 \to K \xrightarrow{\iota} X \xrightarrow{\pi} Q \to 0$ gives a distinguished triangle
$K \to X \to Q \to K[1]$ in $D(\mathcal{A})$: the map $\mathrm{Cone}(\iota) \to Q$, $(a, b) \mapsto \pi(b)$, is a surjective morphism of complexes whose kernel
$K^{n+1} \oplus K^n$ is $\mathrm{Cone}(\id_K)$, which is contractible (proof of
Theorem~\ref{homological-algebra:theorem-homotopy-triangulated}); by the long exact sequence
(Theorem~\ref{homological-algebra:theorem-long-exact-sequence}) it is a quasi-isomorphism, so the triangle of $\iota$ is isomorphic in
$D(\mathcal{A})$ to $K \to X \to Q \to K[1]$. Apply this to $K \to X \to \tau_{\geq n+1}X$ and replace $K$ by $\tau_{\leq n}X$.
\end{proof}

\begin{remark}\label{homological-algebra:remark-t-structures}
The full subcategories $D^{\leq 0}$ and $D^{\geq 0}$ of complexes with cohomology in degrees $\leq 0$, respectively $\geq 0$, form the \emph{standard
$t$-structure} on $D(\mathcal{A})$: $\Hom(D^{\leq 0}, D^{\geq 1}) = 0$, $D^{\leq 0}[1] \subseteq D^{\leq 0}$, and every object sits in a triangle as in
Proposition~\ref{homological-algebra:proposition-truncation}. Its \emph{heart} $D^{\leq 0} \cap D^{\geq 0}$ is equivalent to $\mathcal{A}$. Abstract
$t$-structures, and in particular the perverse $t$-structure on derived categories of sheaves, were introduced in \cite{bbd}; perverse sheaves
and M.~Saito's mixed Hodge modules are the modern setting for the Hodge theory of singular varieties.
\end{remark}

\section{Filtered complexes and strictness}\label{homological-algebra:section-filtered}

\begin{definition}\label{homological-algebra:definition-strict}
A \emph{filtered object} $(A, F)$ is an object with a decreasing filtration $F^pA$, \emph{finite} if $F^aA = A$ and $F^bA = 0$ for some $a, b$. A
morphism $f \colon (A, F) \to (B, F)$ is \emph{filtered} if $f(F^pA) \subseteq F^pB$, and \emph{strict} if $f(F^pA) = f(A) \cap F^pB$ for all $p$. A
morphism of filtered complexes is a \emph{filtered quasi-isomorphism} if it induces quasi-isomorphisms on all $\mathrm{gr}^p$.
\end{definition}

\begin{counterexample}\label{homological-algebra:counterexample-not-strict}
Filtered morphisms need not be strict, and the category of filtered vector spaces is not abelian. Let $A = B = k$ with $F^0A = A$, $F^1A = 0$ and
$F^0B = F^1B = B$, $F^2B = 0$. The identity $A \to B$ is filtered and bijective, but not strict ($f(F^1A) = 0 \neq F^1B \cap B = k$), and it is not an
isomorphism of filtered objects, because its inverse is not filtered. Its kernel and cokernel (with the induced filtrations) are zero, so it is a
monomorphism and an epimorphism; by Lemma~\ref{homological-algebra:lemma-mono-epi-iso} the category is not abelian.
\end{counterexample}

\begin{proposition}[Deligne]\label{homological-algebra:proposition-degeneration-strictness}
Let $K$ be a complex with a filtration that is finite in each degree. The spectral sequence of
Theorem~\ref{homological-algebra:theorem-filtered-complex} degenerates at $E_1$ (that is, $d_r = 0$ for all $r \geq 1$) if and only if every
differential $d \colon K^n \to K^{n+1}$ is strict. In this case $H^n(\mathrm{gr}^pK) \cong \mathrm{gr}^pH^n(K)$.
\end{proposition}

\begin{proof}
We use the notation of Definition~\ref{homological-algebra:definition-z-b}; strictness of $d$ means $F^p \cap d(K) = d(F^p)$.

Suppose $d$ is strict. Let $x \in Z_r^p$ with $r \geq 1$. Then $dx \in F^{p+1} \cap d(K) = d(F^{p+1})$, so $dx = dy$ with $y \in F^{p+1}$. Since
$dy = dx \in F^{p+r}$, $y \in Z_{r-1}^{p+1}$, so $x$ and the cocycle $x - y$ have the same class in $E_r^p$. Hence every class in $E_r^p$ is represented by
a cocycle of $K$, and $d_r = 0$.

Suppose all $d_r$, $r \geq 1$, vanish. Fix a degree $n$. We show $F^p \cap d(K^{n-1}) \subseteq d(F^pK^{n-1})$ for all $p$ by descending
induction on $p$ (for $p$ large, $F^pK^n = 0$). Let $y = dx \in F^p$ with $x \in F^aK^{n-1}$; since the filtration is finite in degree $n - 1$, it suffices
to treat all $a$, by descending induction on $a$. If $a \geq p$ there is nothing to show. If $a < p$, then
$x \in Z_{p-a}^a$, and $d_{p-a}[x] = [dx] = 0$ in $E_{p-a}^p$, so $dx \in B_{p-a}^p + Z_{p-a-1}^{p+1}$: $dx = dx' + z$ with $x' \in F^{a+1}$, $dx' \in F^p$,
and $z \in F^{p+1}$ with $z = d(x - x') \in d(K)$. By induction on $a$ (as $x' \in F^{a+1}$), $dx' \in d(F^p)$; by induction on $p$, $z \in d(F^{p+1}) \subseteq d(F^p)$. So
$y \in d(F^p)$. The last statement follows from $E_1 = E_\infty$ and Theorem~\ref{homological-algebra:theorem-filtered-complex}.
\end{proof}

\begin{remark}\label{homological-algebra:remark-filtered-derived}
Localising the homotopy category of filtered complexes at filtered quasi-isomorphisms gives the \emph{filtered derived category} $DF(\mathcal{A})$.
Deligne used it to construct mixed Hodge structures on the cohomology of algebraic varieties \cite{deligne-hodge-2}; the degeneration of the
Hodge-to-de Rham spectral sequence of a compact K\"ahler manifold (Chapter~\ref{kahler}) is, by
Proposition~\ref{homological-algebra:proposition-degeneration-strictness}, the strictness of the de Rham differential with respect to the
Hodge filtration.
\end{remark}

\section{Exercises}\label{homological-algebra:section-exercises}

\subsection{Homological Algebra}\label{homological-algebra:subsection-exercises-homological-algebra}

\begin{exercise}\label{homological-algebra:exercise-nine-lemma}
Prove the $3 \times 3$ lemma: in a commutative $3 \times 3$ diagram of modules with exact columns, if two of the rows are short exact
sequences, so is the third.
\end{exercise}

\begin{exercise}\label{homological-algebra:exercise-zig-zag}
Let $0 \to A^\bullet \xrightarrow{f} B^\bullet \xrightarrow{g} C^\bullet \to 0$ be a short exact sequence of complexes of modules.
Prove Theorem~\ref{homological-algebra:theorem-long-exact-sequence} directly: show that the recipe of
Remark~\ref{homological-algebra:remark-zig-zag} gives a well defined homomorphism $\delta \colon H^n(C) \to H^{n+1}(A)$, and that
the long sequence is exact at $H^n(B)$, at $H^n(C)$ and at $H^{n+1}(A)$.
\end{exercise}

\begin{solution}
\emph{$a$ is a cocycle:} $f(d(a)) = d(f(a)) = d(d(b)) = 0$ and $f$ is injective. \emph{Independence of $b$:} another preimage
of $c$ is $b + f(a')$ with $a' \in A^n$, since $\ker g = \operatorname{im} f$; it gives $d(b) + f(d(a'))$, hence $a + d(a')$, which has
the same class. \emph{Independence of $c$ in its class:} for $c + d(c')$, write $c' = g(b')$ and use the preimage $b + d(b')$, which
gives the same $d(b)$ and the same $a$. So $\delta$ is well defined, and it is additive because sums of choices are choices for
sums.

\emph{Exactness at $H^n(B)$.} $g_*f_* = (gf)_* = 0$. If $b$ is a cocycle with $g_*[b] = 0$, then $g(b) = d(c')$ with $c' = g(b')$,
so $g(b - d(b')) = 0$ and $b - d(b') = f(a)$ for some $a$; then $f(d(a)) = d(b) = 0$, so $a$ is a cocycle and
$[b] = [b - d(b')] = f_*[a]$.

\emph{Exactness at $H^n(C)$.} If $c = g(b)$ with $b$ a cocycle, use this $b$: $d(b) = 0$, so $\delta[c] = 0$. Conversely, if
$\delta[c] = 0$, then $a = d(a')$ and $b - f(a')$ is a cocycle, since $d(b - f(a')) = f(a) - f(d(a')) = 0$, with
$g(b - f(a')) = c$; so $[c] \in \operatorname{im} g_*$.

\emph{Exactness at $H^{n+1}(A)$.} $f_*\delta[c] = [f(a)] = [d(b)] = 0$. Conversely, if $a$ is a cocycle with $f(a) = d(b)$ for some
$b$, then $c = g(b)$ is a cocycle, since $d(c) = g(d(b)) = g(f(a)) = 0$, and by construction $\delta[c] = [a]$.
\end{solution}

\begin{exercise}\label{homological-algebra:exercise-ext-divisible}
Show that $\Ext^1_\ZZ(\QQ, \ZZ) \neq 0$. (Hint: use $0 \to \ZZ \to \QQ \to \QQ/\ZZ \to 0$.)
\end{exercise}

\begin{exercise}\label{homological-algebra:exercise-tor-torsion-free}
Let $R$ be a PID. Show that an $R$-module is flat if and only if it is torsion-free.
\end{exercise}

\begin{exercise}\label{homological-algebra:exercise-derived-cone}
Let $f \colon A^\bullet \to B^\bullet$ be a morphism of complexes. Show that $f$ is homotopic to zero if and only if the sequence
$0 \to B \to \mathrm{Cone}(f) \to A[1] \to 0$ splits as a sequence of complexes.
\end{exercise}

\begin{exercise}\label{homological-algebra:exercise-injective-dimension}
Show that an abelian group is injective if and only if $\Ext^1_\ZZ(\ZZ/n\ZZ, A) = 0$ for all $n \geq 1$.
\end{exercise}

\begin{solution}
By Example~\ref{homological-algebra:example-ext-tor-cyclic}, $\Ext^1(\ZZ/n\ZZ, A) = A/nA$, which vanishes for all $n$ iff $A$ is
divisible, iff $A$ is injective by Theorem~\ref{homological-algebra:theorem-enough-injectives}.
\end{solution}

\subsection{Spectral Sequences}\label{homological-algebra:subsection-exercises-spectral-sequences}

\begin{exercise}\label{homological-algebra:exercise-two-columns}
State and prove the analogue of Proposition~\ref{homological-algebra:proposition-two-rows} for a spectral sequence concentrated in two
columns $p \in \{0, s\}$.
\end{exercise}

\begin{exercise}\label{homological-algebra:exercise-two-out-of-three}
Let $0 \to A^\bullet \to B^\bullet \to C^\bullet \to 0$ be a short exact sequence of bounded below complexes, viewed as a double complex with three
columns. Using both spectral sequences of Theorem~\ref{homological-algebra:theorem-double-complex} (after shifting so that the double complex is
first-quadrant), show that if two of $A$, $B$, $C$ are acyclic, then so is the third.
\end{exercise}

\begin{exercise}\label{homological-algebra:exercise-balance-ss}
Use Theorem~\ref{homological-algebra:theorem-double-complex} to give a second proof of
Theorem~\ref{homological-algebra:theorem-balance}.
\end{exercise}

\begin{exercise}\label{homological-algebra:exercise-base-change-tor}
Let $R \to S$ be a homomorphism of commutative rings, $M$ an $S$-module and $N$ an $R$-module. Construct a spectral sequence
$\Tor_p^S(M, \Tor_q^R(S, N)) \Rightarrow \Tor_{p+q}^R(M, N)$ (homological indexing).
\end{exercise}

\subsection{Derived Categories}\label{homological-algebra:subsection-exercises-derived-categories}

\begin{exercise}\label{homological-algebra:exercise-five-lemma-triangles}
Show that in a triangulated category, if $(u, v, w)$ is a morphism of distinguished triangles and two of $u$, $v$, $w$ are isomorphisms, so is the
third.
\end{exercise}

\begin{solution}
Let $(u, v, w)$ be a morphism from $X \to Y \to Z \to X[1]$ to $X' \to Y' \to Z' \to X'[1]$ with $u$, $v$ isomorphisms, and let $W$ be any object.
Applying $\Hom(W, -)$ and Proposition~\ref{homological-algebra:proposition-cohomological} gives two long exact sequences of abelian groups
connected by the maps induced by $u$, $v$, $w$, $u[1]$, $v[1]$. The five lemma (Lemma~\ref{homological-algebra:lemma-five}) shows that
$w_* \colon \Hom(W, Z) \to \Hom(W, Z')$ is bijective for all $W$. Taking $W = Z'$ gives $s \colon Z' \to Z$ with $ws = \id$; then $w(sw) = w$, and
injectivity of $w_*$ for $W = Z$ gives $sw = \id$. The other cases follow by rotating the triangles with (TR2).
\end{solution}

\begin{exercise}\label{homological-algebra:exercise-split-complex}
Let $k$ be a field. Show that every complex $X$ of $k$-vector spaces is isomorphic in $D(\cat{Vect}_k)$ to the complex $\bigoplus_n H^n(X)[-n]$ with
zero differential.
\end{exercise}

\begin{solution}
Every subspace has a complement. Choose complements $X^n = Z^n \oplus C^n$ and $Z^n = B^n \oplus H^n$, where $H^n \cong H^n(X)$. The subcomplex
$H = \bigoplus_n H^n$ has zero differential. Choose linear maps $r^n \colon X^n \to H^n$ that are the identity on $H^n$ and vanish on $B^n \oplus C^n$. Then
$r \colon X \to H$ is a morphism of complexes: $r(dx) = 0$ because $dx \in B^{n+1}$. It induces the identity on cohomology, so it is a
quasi-isomorphism, hence an isomorphism in $D(\cat{Vect}_k)$.
\end{solution}

\begin{exercise}\label{homological-algebra:exercise-ext-composition}
Using Proposition~\ref{homological-algebra:proposition-derived-functor-properties}, define the Yoneda product
$\Ext^m(Y, Z) \times \Ext^n(X, Y) \to \Ext^{m+n}(X, Z)$ as composition in $D(\mathcal{A})$ and show that it is associative.
\end{exercise}

\begin{exercise}\label{homological-algebra:exercise-not-abelian}
Show that $D(\mathcal{A})$ is not abelian unless every short exact sequence in $\mathcal{A}$ splits. (Hint: a monomorphism in a triangulated category
is split.)
\end{exercise}

\begin{solution}
Let $f \colon X \to Y$ be a monomorphism in a triangulated category, and complete it to a triangle $Z[-1] \xrightarrow{h} X \xrightarrow{f} Y \to Z$. Then
$fh = 0$, so $h = 0$. Applying Proposition~\ref{homological-algebra:proposition-cohomological} with $\Hom(-, X)$ shows that
$\Hom(Y, X) \to \Hom(X, X)$ is surjective (the next term maps through $h = 0$), so $f$ has a retraction. Dually, epimorphisms have sections.

Suppose $D(\mathcal{A})$ is abelian and let $0 \to A \to B \xrightarrow{g} C \to 0$ be exact in $\mathcal{A}$. As in the proof of
Proposition~\ref{homological-algebra:proposition-truncation}, it gives a distinguished triangle $A \to B \xrightarrow{g} C \xrightarrow{\delta} A[1]$.
Factor $\delta = me$ with $e \colon C \to I$ an epimorphism and $m \colon I \to A[1]$ a monomorphism. Both split, so $I$ is a direct summand of $C$ and of
$A[1]$ in $D(\mathcal{A})$. Cohomology is additive, so $H^k(I)$ vanishes for $k \neq 0$ (as a summand of $C$) and for $k \neq -1$ (as a summand of $A[1]$).
Hence $I$ is acyclic, $I \cong 0$ and $\delta = 0$. By Proposition~\ref{homological-algebra:proposition-cohomological}, the sequence
$\Hom(C, B) \to \Hom(C, C) \to \Hom(C, A[1])$ is exact and $\id_C \mapsto \delta = 0$, so $g\sigma = \id_C$ for some $\sigma \colon C \to B$ in
$D(\mathcal{A})$. Applying $H^0$ gives $g \circ H^0(\sigma) = \id_C$ in $\mathcal{A}$, so the sequence splits.
\end{solution}
